% !TeX program = xelatex
% 第二册第6—8章新手版整合稿
% 主线案例：mini_mec麦轮小车从仿真到实物
% 技术依据：项目x5_src_250401静态审查；ROS 2 Humble实机运行待验证

\documentclass[UTF8,zihao=-4]{ctexart}
\usepackage{amsmath,amssymb,bm}
\usepackage{booktabs}
\usepackage{enumitem}
\usepackage{geometry}
\usepackage{hyperref}
\usepackage{siunitx}
\geometry{a4paper,left=28mm,right=28mm,top=25mm,bottom=25mm}
\hypersetup{colorlinks=true,linkcolor=blue,citecolor=blue,urlcolor=blue}
\sisetup{per-mode=symbol}
\setcounter{tocdepth}{2}
\usepackage{graphicx}
\usepackage{float}
\newcommand{\bookchapter}[1]{%
  \clearpage\part*{#1}\addcontentsline{toc}{part}{#1}}
\newcommand{\booksection}[1]{%
  \section*{#1}\addcontentsline{toc}{section}{#1}}
\newcommand{\booksubsection}[1]{%
  \subsection*{#1}\addcontentsline{toc}{subsection}{#1}}
\DeclareRobustCommand{\code}[1]{\texttt{#1}}
\newcommand{\statusstatic}{\textbf{验证状态：静态审查。}}

\title{《机器人操作系统与智能机器人开发——应用篇》\\
\large 第6—8章：麦轮小车从仿真到现实}
\author{}
\date{}

\begin{document}
\maketitle
\tableofcontents

\section*{贯穿案例：把一辆虚拟小车带到现实世界}

想象一个场景，现在你的面前放着一辆四轮麦克纳姆轮小车。我们希望它完成一个看起来很简单的任务：从起点出发，绕过纸箱，到达教室另一侧的目标点。

如果用键盘遥控，这个任务并不难。但如果拔掉键盘，让小车自己走，问题会一个接一个地出现。

首先，小车必须知道一条速度指令怎样变成四个车轮的转动。接着，它必须知道自己走了多远。然后，它要用激光雷达认识环境并建立地图。最后，它还要在地图上选择路线，并在有人经过时减速、绕行或重新规划。

这正是第6章到第8章要完成的事情。我们不会同时讲完所有概念，而是每遇到一个问题，就引入一个解决问题所必需的工具。

本案例分为两个环境。

\begin{itemize}
  \item \textbf{仿真环境：}Ubuntu 22.04、ROS 2 Humble和与其兼容的Gazebo。仿真机器人应提供\code{/cmd\_vel}、\code{/odom}、\code{/imu/data\_raw}、\code{/scan}和TF。
  \item \textbf{实物环境：}项目\code{x5\_src\_250401}中的\code{mini\_mec}配置。底盘节点接收\code{/cmd\_vel}，发布\code{/odom}和\code{/imu/data\_raw}；激光雷达发布\code{/scan}。
\end{itemize}

这两个环境的内部实现不同，但对上层算法暴露相同接口。这一点很重要。我们不是把仿真程序“改写”成实物程序，而是尽量只替换数据来源和设备驱动。

\begin{center}
任务目标 $\rightarrow$ 状态估计 $\rightarrow$ 建图与定位 $\rightarrow$ 路径规划 $\rightarrow$ 速度控制 $\rightarrow$ 小车运动
\end{center}

\noindent\statusstatic 文中对\code{x5\_src\_250401}的包名、文件路径、话题和主要参数来自本地源码检查。尚未在Ubuntu 22.04、ROS 2 Humble、Gazebo和实物小车上完成运行验证，因此命令和现象均应在教学平台复测后再标记为“已运行验证”。

% ============================================================
\bookchapter{第6章 移动机器人的局部位姿估计}

\section*{本章任务：让小车知道自己刚才怎样运动}

本章先不建图，也不导航。我们只做一件事：让麦轮小车完成前进、横移和旋转，并给出连续的局部位姿。这和小车内部的传感器紧密相关。

学习时请始终盯住三个可观察量：发送给小车的\code{/cmd\_vel}、底盘给出的\code{/odom}，以及融合后的\code{/odom\_combined}。如果这三者之间的关系说不清楚，后面的SLAM和Nav2就没有可靠基础。

\subsection*{学习目标}

完成本章后，读者应能够：

\begin{itemize}
  \item 根据期望底盘速度判断四个麦轮的转动方向；
  \item 解释\code{base\_footprint}、\code{odom}和\code{odom\_combined}的作用；
  \item 区分六轴IMU和九轴IMU实际测量的物理量；
  \item 用\code{robot\_localization}融合轮式里程计与IMU；
  \item 通过直线、横移和原地旋转三个实验定位方向符号、漂移和TF故障。
\end{itemize}

\subsection*{本章主线}

\begin{enumerate}
  \item 在仿真中给出三个最小速度指令。
  \item 用麦轮运动学解释四个车轮为何这样转。
  \item 从编码器和IMU得到局部运动测量。
  \item 用EKF得到连续、较稳定的\code{/odom\_combined}。
  \item 保持接口不变，把同一组检查迁移到\code{mini\_mec}实物小车。
\end{enumerate}

\booksection{6.1 机器人状态与坐标系}


\booksubsection{6.1.1 位姿、速度与运动状态定义}

在平面运动中，小车的位置和朝向合称为位姿。对麦轮小车而言，它通常在一个平面上运动，因此位姿可写成
\begin{equation}
\bm{p}=\begin{bmatrix}x&y&\theta\end{bmatrix}^{\mathrm T},
\end{equation}
其中 $x$ 和 $y$ 表示车体中心在平面坐标系中的位置，$\theta$ 表示车头朝向相对 $x$ 轴的角度。位姿回答的是“小车现在在哪里、朝向哪里”。比如说，若我们的小车初始位于原点、车头朝 $x$ 轴正向，则其位姿为 $[0,0,0]^{\mathrm T}$；若它原地左转 $90^\circ$，则位姿变为 $[0,0,\pi/2]^{\mathrm T}$，位置不变但朝向改变。

接下来看速度。速度向量通常写成
\begin{equation}
\bm{v}=\begin{bmatrix}v_x&v_y&\omega_z\end{bmatrix}^{\mathrm T},
\end{equation}
其中 $v_x$ 和 $v_y$ 是线速度在平面两个轴上的分量，$\omega_z$ 是绕垂直轴旋转的角速度。对麦轮小车来说，它的四个麦克纳姆轮允许车体前后、左右、斜向移动，也可以原地旋转，因此 $v_x$、$v_y$ 和 $\omega_z$ 可以独立控制，不一定满足普通差速车那样的运动约束。

需要注意的是，我们不能直接用速度来代表位置。假设小车以 \SI{0.1}{\metre\per\second} 向前运动，这里的 $0.1\,\text{m/s}$ 只是速度；经过 \SI{5}{\second} 后，小车大约前进 \SI{0.5}{\metre}，这个累计结果才与位置有关。若速度恒定，位置变化可写为
\begin{equation}
\Delta s = v \Delta t = 0.1 \times 5 = 0.5\,\text{m}.
\end{equation}
若速度随时间变化，则需要积分
\begin{equation}
s(t)=s(0)+\int_0^t v(\tau)\,\mathrm{d}\tau.
\end{equation}
在离散估计中，常用近似递推
\begin{equation}
x_{k+1}\approx x_k+v_k\Delta t.
\end{equation}
还要注意，当车体以速度 $v$ 沿自身朝向运动时，速度分量可以表示为
\begin{equation}
v_x=v\cos\theta,\quad v_y=v\sin\theta.
\end{equation}
这说明位姿和速度之间存在非线性耦合：同样的前向速度，在不同朝向下会得到不同的 $x$、$y$ 变化。

最后是状态。状态是算法需要维护的一组量，可以包含位姿、速度、角速度、加速度和传感器偏差。对麦轮小车融合定位问题，一个简单的状态向量可以是
\begin{equation}
\bm{x}=\begin{bmatrix}x&y&\theta&v_x&v_y&\omega_z\end{bmatrix}^{\mathrm T},
\end{equation}
若进一步估计 IMU 或编码器偏差，则可以扩展为
\begin{equation}
\bm{x}=\begin{bmatrix}x&y&\theta&v_x&v_y&\omega_z&b_{\mathrm{gyro}}&b_{\mathrm{enc}}\end{bmatrix}^{\mathrm T}.
\end{equation}
状态不一定都是传感器直接测得的物理量，偏差等量通常需要滤波或优化来估计。

先记住一句话：\textbf{位姿回答“在哪里”，速度回答“正在怎样动”，状态则是估计算法记住的所有相关信息。}

\booksubsection{6.1.2 状态真值、测量值与估计值}

仿真中，我们通常能读取模型的理想位姿，把它近似当作真值。实物中却很难随时知道真正位置。因为真实环境中存在各种噪声和干扰，传感器读数总会与真实值有偏差。所以编码器、IMU和激光雷达只能给出带噪声的测量值。融合算法利用测量值，计算出估计值。

更具体地说，三者含义如下：

\begin{itemize}
  \item \textbf{真值：}指系统在某一时刻的真实状态。仿真中，真值由仿真模型直接给出，例如运动学模型计算出的位姿，通常作为评价基准；但若模型本身存在简化或误差，这个基准与真实物理系统仍有偏差，因此只能称为“参考真值”。实物中，真值无法直接获取，常用高精度动作捕捉系统、全站仪或人工标定值作为近似参考，普通机器人上并不存在绝对真值。
  \item \textbf{测量值：}指传感器直接输出或经过必要换算后得到的数据。例如编码器原始脉冲先换算为轮速或位移，激光雷达原始距离和角度信息被转换为点云坐标。测量值会受到噪声、量化误差、安装偏差和时间延迟等影响，因此是带有不确定性的观测结果。
  \item \textbf{估计值：}指融合算法综合运动模型预测和传感器测量后给出的状态推断结果。估计值不是简单相信模型或简单复制测量，而是通过滤波器或优化器根据不确定性进行加权融合，从而得到比单次测量更可靠的近似。如何计算估计值将会在后面的章节详细赘述。

\end{itemize}

举例来说，某时刻小车真实位姿可能是 $(1.00,0.50,0.20)$；传感器给出的测量可能是 $(1.02,0.48,0.19)$，带有轻微偏差；融合算法输出的估计值可能是 $(1.01,0.49,0.20)$。估计值不必与真值完全相等，但应当稳定且可评价。

如果融合后曲线更平滑，并不能单独证明它更准确。平滑只是外观，准确需要与真值、标定距离或闭环误差比较。过度平滑还可能掩盖真实运动细节，产生延迟或丢失快速转弯。因此，工程上常用绝对轨迹误差、相对位姿误差和终点闭环误差等指标来评价精度，而不是只看曲线是否好看。

另外要注意，估计值不能反过来当作真值去评价另一个估计值，否则容易造成自我验证的假象。

\booksubsection{6.1.3 \code{base\_link}与\code{odom}坐标变换树}

前面提到，位姿 $[x,y,\theta]^{\mathrm T}$ 表示小车在某个参考坐标系中的位置和朝向。但这个参考坐标系并不是唯一的：激光雷达测量障碍物时，距离值是相对于雷达自身中心的；编码器给出的位移是相对于车轮或车体中心的；而导航算法又需要知道小车相对于某个固定世界坐标系的位置。如果大家都用不同的坐标系，数据就无法直接比较和融合。因此，我们必须约定一套坐标系，并明确它们之间的转换关系。

以我们教材的麦轮小车为例，它涉及以下几个常用坐标系：

\begin{itemize}
  \item \textbf{\code{base\_footprint}}：车体在地面上的投影坐标系，原点通常位于车体中心在地面的投影点，随车移动。
  \item \textbf{\code{base\_link}}：车体本体坐标系，原点一般位于车体中心，可能比地面投影高一些，也随车移动。
  \item \textbf{\code{laser}}：激光雷达坐标系，原点位于雷达扫描中心，随雷达安装位置固定，也随车移动。
  \item \textbf{\code{odom\_combined}}：融合里程计所在的局部坐标系，在小车启动后保持局部连续，不随车移动。
\end{itemize}

这些坐标系并不是各自独立的，它们之间有着明确的相对位置关系。例如，激光雷达安装在车体前方 \SI{0.15}{\metre} 处，那么同一个障碍物在雷达坐标系里的坐标和车体坐标系里的坐标就差这 \SI{0.15}{\metre}。这些坐标系之间的关系可以用一棵“变换树”来表示。在 ROS 中，TF（Transform）负责维护这些变换。它保存了任意两个坐标系之间的平移和旋转。把这些关系首尾相接，就形成了一条链：

\begin{center}
\code{odom\_combined} $\rightarrow$ \code{base\_footprint} $\rightarrow$ \code{base\_link} $\rightarrow$ \code{laser}.
\end{center}

这里，\code{odom\_combined} 是相对固定的局部参考坐标系，\code{base\_footprint}、\code{base\_link} 和 \code{laser} 则都随着小车运动。当小车向前移动时，变化的其实是 \code{base\_footprint} 在 \code{odom\_combined} 坐标系中的位置和朝向，而 \code{base\_link} 到 \code{laser} 的变换通常保持不变，因为它由机械安装决定。

为什么不把所有数据都直接放在同一个坐标系里？因为传感器安装在车体不同位置，它们的测量值天然属于各自的坐标系。例如，激光雷达返回前方 \SI{0.5}{\metre} 处有障碍物，这个 \SI{0.5}{\metre} 是相对于雷达中心而言的。如果车体想知道障碍物离车体中心多远，TF 就会利用 \code{laser} 到 \code{base\_link} 的外参，把该点从激光雷达坐标系变换到车体坐标系。这样，导航算法就能在同一参考系下比较障碍物位置和车体位置。

现在来解释“外参”。外参是描述一个坐标系相对于另一个坐标系的固定相对位姿的参数，通常包括平移和旋转两部分。对于平面小车，平移常写成 $[x,y,z]$，表示第二个坐标系原点在第一个坐标系中的位置；旋转可用欧拉角 $[roll,pitch,yaw]$ 或四元数表示，表示两个坐标系之间的朝向差异。例如，激光雷达安装在车体前方 \SI{0.15}{\metre}、高度 \SI{0.30}{\metre} 处，且雷达扫描平面保持水平，那么它相对于车体坐标系的外参可以近似写成 $x=0.15$、$y=0$、$z=0.30$，三个旋转角均为 0。由于传感器与车体之间没有相对运动，安装完成后这些数值不再改变，因此这种变换也被称为“静态变换”。在 TF 中，通常用 \code{static\_transform\_publisher} 来发布这类固定关系。

实物工程中，\code{turn\_on\_wheeltec\_robot.launch.py} 会启动静态变换，并为 \code{mini\_mec} 配置 \code{base\_footprint} 到 \code{laser} 等传感器的外参。这些数值必须与实际硬件安装复核，不能把示例参数当成所有小车的通用尺寸。若外参错误，点云会被投影到错误位置，导致导航和避障判断失效。例如，如果激光雷达的前向偏移量写成了 \SI{0.20}{\metre} 而不是实际的 \SI{0.15}{\metre}，那么雷达探测到的障碍物会被认为离车体更远 \SI{0.05}{\metre}，在窄通道中可能造成碰撞。更隐蔽的是旋转误差：若雷达存在很小的偏航角误差，远处障碍物的横向位置误差会被放大，距离越远偏差越大。

因此，在开始融合定位前，应先检查 TF 树是否完整、外参是否合理。可以用卷尺测量实际安装尺寸，或从 CAD 装配模型中读取设计值，再与启动文件中的参数比对。也可以通过标定方法自动修正外参，进一步提高传感器数据与车体坐标系的对齐精度。
\booksection{6.2 轮式里程计与IMU测量}

在 6.1 节中我们区分了位姿、速度和状态。我们强调了一件事：传感器只能提供带有噪声的测量值，估计器需要综合模型和测量才能得到可靠的状态估计。本章介绍两种最基础的传感器——轮式编码器和惯性测量单元（IMU）。它们分别提供关于车轮转动和车体惯性运动的测量，是后续融合定位算法的重要输入。我们将从最底层的脉冲计数开始，逐步推导出可用于状态估计的速度和位姿增量，并讨论两种传感器各自的误差特性。

\booksubsection{6.2.1 编码器脉冲转轮速及里程计积分}

里程计（Odometry）是一种利用运动传感器数据推算机器人位姿变化的方法。在轮式机器人上，最常用的运动传感器就是编码器。编码器安装在电机轴或车轮上，负责把转动转换为电脉冲；而里程计则负责把这些脉冲“翻译”成车轮转过的角度、车体的速度，再通过积分得到位姿增量。简单地说，编码器提供原始测量，里程计完成从测量到位姿的推算。

编码器不会直接说“小车前进了\SI{0.5}{\metre}”。它只记录电机轴或车轮转过的脉冲数。要得到线速度或角速度，必须经过换算。这个换算过程通常由底盘驱动节点完成，但理解其原理对后续调试至关重要。

假设某编码器每转一圈产生 $N$ 个有效计数（例如 $N=1024$ 或 $N=4096$，具体取决于编码器类型和倍频方式）。在采样间隔 $\Delta t$ 内，第 $i$ 个车轮的编码器计数增加了 $\Delta n_i$，则该车轮的转角增量近似为
\begin{equation}
\Delta\phi_i=2\pi\frac{\Delta n_i}{N}.
\end{equation}
角速度近似为
\begin{equation}
\omega_i=\frac{\Delta\phi_i}{\Delta t}.
\end{equation}
这里用“近似”是因为计数是离散的，且忽略了采样间隔内转速的微小波动。当 $\Delta t$ 足够小时，这种近似可以接受。

对于麦轮小车，四个轮子的角速度与车体速度之间的关系由麦克纳姆轮正运动学决定。记四个轮子的角速度为 $\omega_1,\omega_2,\omega_3,\omega_4$，车体速度为 $[v_x,v_y,\omega_z]^{\mathrm T}$，其中 $v_x$、$v_y$ 是车体坐标系下的线速度，$\omega_z$ 是绕车体中心的旋转角速度。根据第5章推导的正运动学公式，可以写出
\begin{equation}
\begin{bmatrix} v_x \\ v_y \\ \omega_z \end{bmatrix}
= \bm{J}
\begin{bmatrix} \omega_1 \\ \omega_2 \\ \omega_3 \\ \omega_4 \end{bmatrix},
\end{equation}
其中矩阵 $\bm{J}$ 由轮子半径、辊子安装角以及车体几何尺寸决定。具体表达式见第5章，此处不再重复推导。重要的是：一旦得到车体速度，就可以通过积分得到位姿变化。

在离散时间下，假设 $k$ 时刻的位姿为 $[x_k,y_k,\theta_k]^{\mathrm T}$，经过一个采样间隔 $\Delta t$，新的位姿可近似为
\begin{align}
x_{k+1}&=x_k+(v_x\cos\theta_k-v_y\sin\theta_k)\Delta t,\\
y_{k+1}&=y_k+(v_x\sin\theta_k+v_y\cos\theta_k)\Delta t,\\
\theta_{k+1}&=\theta_k+\omega_z\Delta t.
\end{align}
前两个式子把车体坐标系下的速度旋转到世界坐标系：车体前进方向的速度 $v_x$ 在世界坐标系下会同时贡献 $x$ 和 $y$ 的变化，具体取决于当前朝向 $\theta_k$。第三个式子说明朝向角的变化等于角速度乘以时间。如果 $\omega_z$ 也是时间的函数，严格地说应当积分 $\theta(t)=\theta(0)+\int_0^t \omega_z(\tau)\,d\tau$；在离散实现中，用当前采样值乘以 $\Delta t$ 是一种欧拉近似。

这里最重要的不是记住三行公式，而是看见“累计”。每一时刻的速度误差都会在下一时刻被带入位姿，而位姿误差又会继续影响后续的速度分解，形成累积漂移。例如，若编码器给出的轮速始终偏大 $1\%$，则 $v_x$ 偏大 $1\%$，在 $\Delta t=0.1$ s 内每个周期的位移误差可能很小，但经过一分钟、两分钟后，累计的距离误差就会达到分米甚至米级。因此，里程计适合描述短时间连续运动，却会随着路程增长而漂移。

在\code{x5\_src\_250401}中，\code{turn\_on\_wheeltec\_robot}的底盘节点接收\code{cmd\_vel}，并发布\code{nav\_msgs/msg/Odometry}类型的\code{odom}。该消息中通常包含以下关键字段：
\begin{itemize}
  \item \code{pose.pose}：根据里程计积分得到的机器人位姿；
  \item \code{twist.twist}：机器人当前速度（车体坐标系下）；
  \item \code{header.frame\_id}：通常为 \code{odom}，表示该位姿是相对于哪个坐标系给出的；
  \item \code{child\_frame\_id}：通常为 \code{base\_footprint} 或 \code{base\_link}，表示该位姿描述的是哪个坐标系。
\end{itemize}


\booksubsection{6.2.2 IMU角速度、加速度与姿态测量}

上一节中，里程计依靠车轮编码器间接推算车体运动。但轮子打滑、空转或麦轮滚子切换时，这种推算就会出错。我们还需要一种直接测量车体自身旋转和加速度的传感器，这就是惯性测量单元（Inertial Measurement Unit，IMU）。

IMU 安装在车体上，随车一起运动，内部通常包含多个微机电（MEMS）传感器，用来感知车体的角速度和比力。常用的IMU分为六轴和九轴两种，这里的“轴”表示测量通道。
\begin{itemize}
  \item \textbf{六轴IMU：}三轴陀螺仪加三轴加速度计。它直接测量 $x,y,z$ 三个方向的角速度，以及 $x,y,z$ 三个方向的比力。静止时，加速度计读数通常包含重力影响。这里的“比力”是指单位质量受到的除重力以外的接触力，因此加速度计测量值 $\bm{a}$ 并不等于实际加速度 $\ddot{\bm{p}}$，而是满足
  \begin{equation}
    \bm{a} = \ddot{\bm{p}} - \bm{g},
  \end{equation}
  其中 $\bm{g}$ 是重力加速度向量。当机器人水平静止时，$\ddot{\bm{p}}=0$，加速度计读数应接近 $[0,0,g]^{\mathrm T}$（具体符号取决于坐标定义）。
  \item \textbf{九轴IMU：}在六轴基础上增加三轴磁力计。磁力计测量三个方向的磁场分量，可为航向提供绝对参考，但容易受电机、电流和钢铁结构干扰。磁力计的读数单位通常是微特斯拉（$\mu$T），其方向指向地磁北极附近，与地理北极存在磁偏角。
\end{itemize}

无论六轴还是九轴，IMU 直接输出的都是角速度、比力或磁场等原始物理量，而不是现成的姿态角。要得到稳定可靠的姿态估计，还需要通过滤波或优化算法把这些原始测量融合起来。因此，在使用 IMU 数据时，必须清楚地区分“传感器直接测量”和“算法估计结果”这两个层次。得到了原始物理量之后，还需要通过算法估计才能得到姿态角。


\booksubsection{6.2.3 里程计累计误差与IMU零偏、漂移特性}

轮式里程计相信车轮。轮子空转时，它以为车动了；小车被外力推走但车轮没有正确转动时，它又可能低估运动。具体来说，里程计误差主要来源于：
\begin{itemize}
  \item \textbf{轮径误差：}实际轮子直径与标称值不符，或轮胎磨损、气压变化导致有效半径改变；
  \item \textbf{轮距和轴距误差：}车体几何参数不准确，使旋转运动解算产生系统偏差；
  \item \textbf{打滑和空转：}轮子与地面之间出现相对滑动，编码器计数与实际位移不再一致；
  \item \textbf{麦轮滚子切换：}麦克纳姆轮的小滚子在不同接触阶段会产生非线性效应，尤其在高速或横移时，$y$ 方向速度估计可能产生周期性误差；
  \item \textbf{采样和量化：}编码器分辨率有限，低速时计数更新缓慢，速度估计噪声增大。
\end{itemize}

IMU相信惯性测量。陀螺仪存在零偏，小偏差积分后会变成不断增长的角度误差；加速度经过两次积分才得到位置，误差增长更快；磁力计则可能被周围磁场扰动。具体来说，IMU误差主要来源于：
\begin{itemize}
  \item \textbf{陀螺仪零偏：}即使静止，角速度输出也可能不为零。这个零偏会随时间缓慢变化，积分后导致姿态持续漂移；
  \item \textbf{加速度计零偏和噪声：}加速度计的偏置和随机噪声会在一次积分中造成速度误差，两次积分后造成位置误差，且随时间平方增长；
  \item \textbf{振动敏感性：}电机运转、路面不平引起的振动会耦合进加速度计，产生高频噪声，有时还会使陀螺仪短时饱和；
  \item \textbf{磁力计干扰：}电机、电池、金属结构产生的软磁和硬磁干扰会扭曲磁场读数，使航向估计偏离真实方向。
\end{itemize}

因此，编码器和IMU不是“谁替代谁”的关系。编码器提供平移运动的重要线索，陀螺仪能更直接地观察旋转。二者组合通常比单独使用更稳健。例如，里程计可以提供短时间内的相对位移，而陀螺仪可以校正航向的快速变化；加速度计和磁力计则可以在长时间尺度上辅助修正姿态漂移。融合算法的任务就是根据各自的误差特性动态加权，给出一个比任何单一传感器都更可靠的位姿估计。这为第6.3节介绍的融合定位方法奠定了传感器基础。

\booksection{6.3 多传感器融合（EKF）}

前面两节分别介绍了轮式里程计和 IMU 的测量原理与误差特性。轮式里程计在短时间内能提供平滑的位姿增量，但会随着路程累积漂移；IMU 能直接感知旋转和比力，但存在零偏和噪声，积分后误差增长迅速。单独使用任何一个传感器都难以长时间保持可靠的位姿估计。本章介绍的扩展卡尔曼滤波（Extended Kalman Filter，EKF）将多个传感器的测量与运动模型结合起来，利用各自的优势互补，从而得到比单一传感器更稳健的估计结果。

\booksubsection{6.3.1 为什么融合：互补滤波思想}

假设小车收到“保持当前速度”的指令。即使下一帧传感器数据稍晚到达，我们也可以先根据上一时刻速度猜测新位置。这叫预测。

随后，编码器和 IMU 送来新测量。若测量可靠，就把估计向测量方向拉近；若测量噪声大，就少相信一些。这叫更新或纠正。

这个“预测—更新”的循环正是卡尔曼滤波的核心。扩展卡尔曼滤波（Extended Kalman Filter，EKF）是卡尔曼滤波在非线性系统上的推广，其核心直觉仍然是：

\begin{center}
旧估计 $+$ 运动模型 $\rightarrow$ 预测；\qquad 预测 $+$ 新测量 $\rightarrow$ 更新。
\end{center}

预测步骤利用机器人的运动模型，根据上一时刻的状态和控制输入（如速度指令）推算当前时刻的状态。这个预测值并不是凭空猜测，而是基于物理规律的先验估计。但由于模型误差和未建模因素，预测值会带有不确定性。

更新步骤利用传感器测量值来修正预测。传感器测量同样带有噪声，因此不能完全相信测量，也不能完全相信预测。卡尔曼滤波根据两者各自的不确定性进行加权融合：不确定性小的信息获得更高的权重。这就是“互补滤波”思想的体现——不同传感器在不同频率范围内各有优势，融合算法让它们在各自擅长的区间发挥作用。

协方差表示“不确定到什么程度”。它不是随便填的权重。协方差越小，滤波器越信任该测量；若把所有协方差都设得很小，不同传感器互相冲突时，结果可能抖动甚至跳变。例如，若把 IMU 的角速度测量协方差设得过小，而陀螺仪存在零偏，滤波器会过度相信有偏的角速度，导致航向估计被持续拉偏。反之，协方差过大则会使滤波器对测量反应迟钝，无法及时纠正误差。因此，协方差的设置需要结合传感器实际噪声水平，不能随意填写。

\booksubsection{6.3.2 扩展卡尔曼滤波入门：预测与更新}

EKF 将状态估计问题分为预测和更新两步。以平面移动小车为例，状态向量通常包含位置、朝向和速度：
\begin{equation}
\bm{x} = \begin{bmatrix} x & y & \theta & v_x & v_y & \omega_z \end{bmatrix}^{\mathrm T}.
\end{equation}
预测阶段利用上一时刻的状态和运动模型，例如上一小节提到的离散递推关系：
\begin{align}
x_{k+1}&=x_k+(v_x\cos\theta_k-v_y\sin\theta_k)\Delta t,\\
y_{k+1}&=y_k+(v_x\sin\theta_k+v_y\cos\theta_k)\Delta t,\\
\theta_{k+1}&=\theta_k+\omega_z\Delta t.
\end{align}
同时，预测阶段还会根据过程噪声更新协方差矩阵，表示状态不确定性的增加。更新阶段则将编码器、IMU 等传感器的测量值融入估计，校正预测偏差，并降低协方差。EKF 通过线性化非线性模型来处理这些关系，具体公式在机器人学教材中都有详细推导，这里不展开。

实物工程通过 \code{turn\_on\_wheeltec\_robot/launch/wheeltec\_ekf.launch.py} 启动 \code{robot\_localization} 中的 \code{ekf\_node}。配置文件为 \code{config/ekf.yaml}。\code{robot\_localization} 是一个通用的多传感器融合包，内部实现了 EKF 和无迹卡尔曼滤波（UKF），用户只需通过 YAML 文件配置输入话题、状态量映射和协方差即可，无需修改核心算法。

该配置的关键数据流是：

\begin{center}
\code{/odom} $+$ \code{/imu/data\_raw} $\rightarrow$ \code{ekf\_filter\_node} $\rightarrow$ \code{/odom\_combined}.
\end{center}

这里 \code{/odom} 是底盘节点发布的轮式里程计消息，提供位姿和速度；\code{/imu/data\_raw} 是 IMU 原始数据，提供角速度和加速度；\code{ekf\_filter\_node} 融合两者后，输出 \code{/odom\_combined}，即融合后的里程计消息。通过 TF，\code{/odom\_combined} 成为 \code{base\_footprint} 的参考坐标系，后续导航和可视化都基于它。

配置设置 \code{two\_d\_mode: true}，表示只估计平面运动，状态向量退化为 $[x,y,\theta,v_x,v_y,\omega_z]^{\mathrm T}$，忽略 $z$、横滚和俯仰。\code{world\_frame} 和 \code{odom\_frame} 使用 \code{odom\_combined}，\code{base\_link\_frame} 使用 \code{base\_footprint}。这表示融合结果被发布在 \code{odom\_combined} 坐标系下，描述 \code{base\_footprint} 相对于该坐标系的位姿。

阅读长 YAML 文件时，不要从第一行读到最后一行。先回答四个问题：输入话题是什么、每个输入启用了哪些状态量、输出坐标系是什么、是否发布 TF。找到这四项，配置就有了骨架。例如，在 \code{ekf.yaml} 中通常会有 \code{odom0: /odom}、\code{imu0: /imu/data\_raw} 等输入定义，每个输入下方列出了 \code{pose}、\code{twist}、\code{imu} 等字段及启用标记；输出由 \code{odom\_frame}、\code{base\_link\_frame} 等参数指定；\code{publish\_tf: true} 表示该节点会发布 \code{odom\_combined} 到 \code{base\_footprint} 的 TF。抓住这些关键点，配置结构就清晰了。

\booksubsection{6.3.3 \code{robot\_localization}融合配置与协方差调优}

调参从静止开始，而不是从复杂导航开始。原因很简单：静止时没有任何外部运动干扰，传感器读数应只反映噪声和零偏，便于观察滤波器的初始化行为和漂移趋势。如果静止状态下融合位姿都出现明显漂移，说明基础配置有问题，应先解决再考虑动态调参。

建议按以下步骤进行：

\begin{enumerate}
  \item 静止记录 \code{/odom}、\code{/imu/data\_raw} 和 \code{/odom\_combined}。使用 \code{ros2 bag} 或 \code{rqt\_plot} 观察数据。
  \item 检查静止时融合位姿是否缓慢旋转或平移。例如，航向角是否持续增大，位置是否缓慢漂移。若航向漂移明显，可能与陀螺仪零偏、时间戳对齐、坐标方向定义或重力处理有关。
  \item 让小车低速直行一段已知距离，比较原始与融合结果。可以在起点和终点做标记，记录 \code{/odom} 和 \code{/odom\_combined} 给出的位移量，与真实距离对比。
  \item 再做横移和原地旋转，每次只改变一类参数。例如先调整与 IMU 相关的协方差，观察航向漂移变化；再调整与里程计相关的协方差，观察平移精度变化。避免同时修改多个参数，否则难以判断哪个参数起了作用。
\end{enumerate}

若静止时航向持续漂移，先检查 IMU 零偏、时间戳、坐标方向和重力处理，再考虑增大相应测量协方差。不要用极端协方差把坏数据“藏起来”。比如，把陀螺仪测量协方差调到非常大，确实会抑制漂移，但同时也会让滤波器几乎忽略 IMU 的角速度测量，退化回仅靠里程计的状态，失去了融合的意义。正确的做法是标定 IMU 零偏，确保时间戳同步，统一坐标系方向，然后设置合理的协方差。

协方差的初始值可以参考传感器数据手册中的噪声密度，再结合实测方差进行调整。例如，让小车静止，记录 \code{/imu/data\_raw} 中角速度的标准差，作为陀螺仪测量噪声协方差的参考；记录编码器推算速度的波动，作为里程计速度协方差的参考。这些实测值比盲目猜测更有依据。

调参是一个迭代过程，需要耐心和记录。每次修改配置后，重新运行测试，记录结果，逐步逼近满意的精度。注意，融合位姿的平滑并不等于准确，最终仍要结合真实位移、闭环误差或外部参考进行评价。
\booksection{6.4 传感器数据预处理}

\booksubsection{6.4.1 消息时间戳对齐与频率匹配}

在融合算法中，时间同步是一个容易被忽略却至关重要的问题。融合算法本质上是把所有传感器数据放到同一条时间线上进行比较和组合，它默认每个输入数据都描述同一时刻的机器人状态。如果编码器消息的时间戳写着“现在”，但 IMU 数据实际上来自 0.2 秒之前，滤波器就会把两个不同时刻的运动状态强行拼在一起，就像把两张不同时间拍的照片叠起来找差异，结果必然产生畸变。因此，检查消息头中的 header.stamp 是必要的。理想情况下，时间戳应尽量接近传感器的实际采样时刻，而不是消息被发布或接收的时刻。有些驱动节点会在采样瞬间打上时间戳，但有些设备或驱动可能使用接收时刻，这会引入固定或可变延迟，需要留意。

除了时间戳，header.frame\_id 同样重要。它指明数据所在的坐标系。融合节点需要把不同传感器的数据变换到统一坐标系下才能比较，如果某个消息的 frame\_id 在 TF 树中不存在，或者名称写错，坐标变换就会失败，融合结果也会出错。因此，在配置融合之前，应确认每个输入话题的 frame\_id 都能在 TF 树中找到，并且与实际的物理安装一致。

不同传感器的发布频率不必完全相同。例如，IMU 可能以 100 Hz 甚至更高频率发布，而轮式里程计可能只有 20 Hz 或 50 Hz。滤波器会根据时间戳对数据进行插值或等待，因此频率不同并不妨碍融合。但频率应当稳定，忽快忽慢或长时间无数据会导致预测步骤长时间得不到校正，误差迅速积累。真正危险的情况包括：某个传感器长时间不更新，滤波器只能依靠模型预测，漂移会逐渐增大；时间戳倒退（例如时钟跳变或驱动错误）会破坏滤波器的因果假设；延迟突然增大意味着测量值严重滞后于真实状态，融合结果会出现震荡或偏差。因此，在调试过程中，不仅要关注数据内容，还要检查消息的时间戳和发布频率是否正常。


\booksubsection{6.4.2 异常值、丢包与数据超时处理}

在真实系统中，传感器数据并不总是干净可靠的。电磁干扰、通信抖动、驱动错误或物理冲击都可能产生异常测量。融合算法如果直接使用这些数据，轻则估计结果抖动，重则发散。因此，在数据进入滤波器之前，需要设置一层简单的规则，把明显错误的值挡在门外。

常见的异常包括：
\begin{itemize}
  \item \textbf{非数值（NaN）和无限大（Inf）}：可能由传感器故障或除零导致，必须直接丢弃，否则会使滤波器输出无效状态。
  \item \textbf{超过物理上限的速度}：例如麦轮小车设计最大速度为 \SI{1.5}{\metre\per\second}，若编码器给出的速度突然变成 \SI{10}{\metre\per\second}，显然不合理，可以限幅或丢弃。
  \item \textbf{异常跳变的角速度}：IMU 角速度在短时间内不应突变，若连续两帧差值超过阈值（如大于 \SI{3}{\radian\per\second} 的跳变），可能是数据尖峰或通信错误。
  \item \textbf{时间戳倒退}：消息的时间戳不应早于上一帧，否则会破坏滤波器的时序假设，需要丢弃并记录警告。
\end{itemize}

处理策略必须明确：是直接丢弃该帧，还是用上一帧值替换，还是进行限幅。例如，对于速度超限，可以将其限制在最大允许值；对于 NaN 或 Inf，只能丢弃；对于时间戳倒退，通常丢弃该帧并重新同步。不要简单地把所有异常都交给“滤波”处理，因为滤波器本身可能被极端异常值带偏。

此外，需要设置连续异常计数。如果偶尔出现一两个异常点，可以忽略；但如果连续很多帧异常（例如连续 10 次或 100 ms 内持续异常），说明传感器可能已经故障，系统应当报警，并根据应用需求决定是否进入安全停止或切换备用传感器。可以在代码中维护一个计数器，一旦出现异常就累加，收到正常数据就清零；当计数器超过阈值时触发保护逻辑。

对于丢包，可以使用 \code{sensor\_timeout} 参数设定一个超时时间。如果某个传感器在超时时间内没有新数据到来，滤波器只进行预测而不更新，同时增大协方差以反映不确定性增加。如果超时时间过长，预测误差会累积，因此需要合理设置。一般建议将超时时间设置为传感器正常发布周期的 2~3 倍。例如，里程计以 \SI{20}{\hertz} 发布，则 \code{sensor\_timeout} 可设为 \SI{0.1}{\second} 到 \SI{0.15}{\second}。注意，超时后恢复数据时，滤波器会自动重新开始更新，但需要一个短暂的收敛过程。

通过上述规则，可以过滤掉大部分明显错误，保护融合算法的稳定性。但需要注意，这些规则只是第一道防线，更复杂的异常检测（如基于模型残差的一致性检验）可能需要进一步处理。

\booksubsection{6.4.3 ROS Odometry 消息结构与发布接口}

在 ROS 2 中，里程计数据使用 \code{nav\_msgs/msg/Odometry} 消息类型，它是移动机器人中最常用的消息之一。该消息包含两组核心信息：\code{pose} 和 \code{twist}。

\begin{itemize}
  \item \textbf{pose}：表示机器人相对于父坐标系的位置和姿态，类型为 \code{geometry\_msgs/PoseWithCovariance}。其中 \code{position} 包含 $x,y,z$，\code{orientation} 包含四元数。协方差矩阵描述位置和姿态估计的不确定性，如果底盘节点没有提供协方差，融合算法通常会忽略该测量或赋予一个较大的默认协方差。
  \item \textbf{twist}：表示机器人在子坐标系语境下的线速度和角速度，类型为 \code{geometry\_msgs/TwistWithCovariance}。其中 \code{linear} 包含 $x,y,z$ 的线速度，\code{angular} 包含 $x,y,z$ 的角速度。协方差矩阵描述速度估计的不确定性。
\end{itemize}

同时，消息头中的 \code{header.frame\_id} 和 \code{child\_frame\_id} 定义了坐标系关系。\code{header.frame\_id} 通常设置为父坐标系（参考系），如 \code{"odom"} 或 \code{"odom\_combined"}；\code{child\_frame\_id} 设置为子坐标系（本体坐标系），如 \code{"base\_footprint"} 或 \code{"base\_link"}。\code{pose} 字段描述的是 \code{child\_frame\_id} 在 \code{header.frame\_id} 中的位姿。\code{twist} 字段则通常表示在 \code{child\_frame\_id} 下表达的速度（即机器人本体坐标系下的速度），也有部分驱动在父坐标系下表达，因此必须阅读具体文档或源码确认。

新手最常见的错误是只关注 $x$ 和 $y$ 的数值，而忽略了坐标系。例如，把里程计的 $x$ 直接当作地图坐标系的 $x$，或者把不同坐标系的速度直接比较。必须记住：任何位姿和速度都是相对于某个参考系的，脱离坐标系谈数值是没有意义的。在调试时，应养成先看 \code{frame\_id} 的习惯。

检查消息时，可以用 \code{ros2 topic echo /odom} 查看 \code{header.frame\_id} 和 \code{child\_frame\_id}，确认它们与 TF 树中的坐标系一致。如果 \code{frame\_id} 写错或不存在，下游节点无法进行坐标变换，可能导致数据被忽略或错误处理。

底盘节点在发布 \code{Odometry} 消息时，应确保时间戳为当前采样时刻，并且发布频率稳定。同时，要注意 \code{pose} 和 \code{twist} 的协方差是否合理，因为融合算法会根据协方差来加权。如果底盘发布的协方差全为零或未设置，\code{robot\_localization} 等融合节点可能会拒绝使用或默认给一个较大的值，导致融合效果下降。

总之，理解 \code{Odometry} 消息的结构和坐标系约定，是调试和配置融合系统的基础。只有正确解析消息内容，才能确保后续的融合算法获得可靠的输入。

\booksection{6.5 本章验证与排故}

\booksubsection{6.5.1 Gazebo仿真与实物直线、横移测试}

迁移前，把仿真和实物逐项对齐。

\begin{center}
\begin{tabular}{lll}
\toprule
功能 & 仿真接口 & 实物接口 \\
\midrule
速度命令 & \code{/cmd\_vel} & \code{/cmd\_vel} \\
原始里程计 & \code{/odom} & \code{/odom} \\
IMU & \code{/imu/data\_raw} & \code{/imu/data\_raw} \\
融合里程计 & \code{/odom\_combined} & \code{/odom\_combined} \\
激光扫描 & \code{/scan} & \code{/scan} \\
车体坐标系 & \code{base\_footprint} & \code{base\_footprint} \\
\bottomrule
\end{tabular}
\end{center}

接口名称相同并不代表数据一定正确。还要核对单位、方向、发布频率、时间来源、协方差和TF父子关系。

\booksubsection{6.5.2 原始里程计与融合结果比较}

\begin{verbatim}
# ROS 2 Humble；实物mini_mec；来自本地源码，待运行验证
ros2 launch turn_on_wheeltec_robot turn_on_wheeltec_robot.launch.py
ros2 topic list
ros2 topic hz /odom
ros2 topic hz /imu/data_raw
ros2 topic echo /odom_combined --once
\end{verbatim}

预期可以看到底盘、IMU、EKF和机器人模型相关节点与话题。若串口打不开，先检查设备节点、权限和端口配置；若没有\code{/odom\_combined}，再检查EKF节点是否启动以及输入是否到达。

实物第一次运动必须低速、架空或留出安全空间，并准备立即发送零速度。不要从导航故障开始排查一辆尚未验证基本轮向的小车。

\booksubsection{6.5.3 典型故障：静止漂移、横向方向错误与TF跳变}

\begin{enumerate}
  \item \textbf{直行：}给定低速$v_x$，检查$x$主要变化，$y$和航向变化较小。
  \item \textbf{横移：}给定低速$v_y$，检查$y$方向、轮序和麦轮滚子安装。
  \item \textbf{旋转：}给定低速$\omega_z$，检查IMU角速度和里程计航向符号一致。
\end{enumerate}

这三个实验通过后，才进入第7章。否则，SLAM可能把底层方向错误表现成地图重影，让问题更难定位。

\booksection{6.6 本章小结与习题}

本章让小车获得了“短时间内知道自己怎样运动”的能力。\code{/cmd\_vel}描述希望怎样动，编码器和IMU描述实际怎样动，EKF输出连续的\code{/odom\_combined}。这还不是全球位置，但它是SLAM和导航的地基。

\booksubsection{6.6.1 麦克纳姆轮运动学计算题}
\begin{enumerate}
  \item 用一句话分别解释位姿、速度、测量值和估计值。
  \item 当$v_x=0$、$v_y>0$、$\omega_z=0$时，根据矩阵判断四轮组合。
  \item 说明六轴IMU比九轴IMU少了什么，以及新增传感器主要提供哪类参考。
\end{enumerate}

\booksubsection{6.6.2 里程计和IMU融合任务}
记录直行、横移、旋转各\SI{10}{\second}时的\code{/odom}和\code{/odom\_combined}。提交话题频率、终点位姿和你认为最明显的一类误差。

\booksubsection{6.6.3 状态估计故障诊断题}
故意把仿真中的横向速度符号反转。描述现象，并给出从\code{/cmd\_vel}、轮速、\code{/odom}到TF的检查顺序。

\paragraph{Technical English}
尝试用英文解释：``Odometry is locally continuous, but its error accumulates over time.''指出\emph{locally continuous}和\emph{accumulates}分别对应本章哪个现象。

% ============================================================
\bookchapter{第7章 激光SLAM与地图构建}

\section*{本章任务：让小车画出房间，并在地图中找到自己}

第6章我们的小车学会了两种报告自己运动的方式：一是靠轮子转动的圈数来推算走了多远（轮式里程计），二是靠IMU（惯性测量单元）来感知加速度和旋转速度。这两种方式单独使用的时候，都有一个共同的毛病——\textbf{走得越久，误差越大}。轮子会打滑，IMU会积分漂移，就像闭着眼睛走路，数着步子，但走着走着就偏了方向。

现在，我们给小车装上一只眼睛——二维激光雷达。激光雷达不像轮子或IMU那样只能感知自身的运动，它能直接“看见”周围的环境：墙壁在哪里、柱子有多远、桌角在什么方向。有了这些外部参照物，小车就能把自己看到的东西和之前的记忆进行比对，用环境中固定不变的物体来修正自己估计的位置，同时一点一点地把整个房间画成一张地图。

本章仍然使用同一辆麦轮小车和同一组接口。新增的关键输入是\code{/scan}（激光雷达数据），新增的全局坐标系是\code{map}（地图坐标系）。

\booksection{7.1 SLAM任务与坐标系统}

\booksubsection{7.1.1 同时定位与建图问题定义}

同时定位与建图（Simultaneous Localization and Mapping，SLAM）——这个名字听起来很唬人，但其实它的核心问题特别简单：

\begin{quote}
\emph{地图未知时，机器人要估计自己的位置；同时，又得用自己估计的位置去建地图。}
\end{quote}

这就陷入了一个“先有鸡还是先有蛋”的困境。为什么不能先定位再建图？因为定位需要一张地图做参考。为什么不能先建图再定位？因为每束激光必须知道自己是从哪里发射出去的，不知道位置就没法把激光点放到正确的位置上。

两者互相依赖，所以算法只能两条腿走路——交替改进位置估计和地图。这就是SLAM的精髓。

\begin{center}
\fbox{
\begin{minipage}{0.85\textwidth}
\small
\textbf{用生活中的例子来理解}

想象你被蒙着眼带进一个陌生的房间，然后摘下眼罩。你手里拿着一支笔和一张纸，任务是画出房间的俯视图，同时还要在图上标出你自己站在哪里。

你唯一能做的事情就是：伸出手摸一摸周围的墙壁和家具。每摸到一个物体，你就能知道它相对于你的方向和距离。

你开始摸——左边有一面墙，距离2米；前方有一个桌子，距离1.5米……你把这些信息记在纸上。然后你往前走了一步，再摸——咦，前方桌子的距离变成了1米。你知道自己往前走了1米，所以把纸上的自己位置也往前移了1米。

你继续走，继续摸，继续画。走了一会儿，你突然发现：前方摸到的墙壁，和最开始摸到的那面墙好像是同一面！这时候你就知道——哦，我绕了一圈回来了。于是你把纸上画的那面墙“连起来”，形成一个完整的房间轮廓。

这就是SLAM——边走边摸，边摸边画，画着画着发现自己走回了老地方，于是把图画得更准确。
\end{minipage}
}
\end{center}

把上面的故事翻译成数学语言，就变成了这样：

设机器人位姿序列为 \(X = \{x_1, x_2, \dots, x_T\}\)，环境地图为 \(M\)，激光观测为 \(Z = \{z_1, z_2, \dots, z_T\}\)，里程计输入为 \(U = \{u_1, u_2, \dots, u_T\}\)。SLAM要做的就是找到最可能的位姿和地图：

\[
X^*, M^* = \arg\max_{X, M} \, P(X, M \mid Z, U)
\]

这个公式的意思很简单：\textbf{在所有可能的轨迹和地图中，挑出最能解释我们观测到的那一组}。

这个联合分布可以分解为\textbf{运动模型} \(P(x_t \mid x_{t-1}, u_t)\)（根据上一时刻位置和里程计，推测当前在哪里）和\textbf{观测模型} \(P(z_t \mid x_t, M)\)（给定位置和地图，看到某束激光的概率）的乘积。

在实际的激光SLAM中，观测模型的核心是\textbf{扫描匹配}（Scan Matching）——把当前激光扫描和已有的地图（或之前的扫描）进行对齐，算出机器人相对移动了多少。这些匹配结果和里程计数据一起，构成了一个\textbf{位姿图}（Pose Graph），然后通过\textbf{图优化}（Graph Optimization）来求解全局一致的轨迹和地图。

\begin{center}
\fbox{
\begin{minipage}{0.85\textwidth}
\small
\textbf{对我们的小车来说，这意味着什么？}

麦轮小车同时拥有轮式里程计和IMU，它们都在\textbf{递推}地估计位置——每一小步都依赖前一步的结果，所以误差会逐渐累积。激光雷达则提供了相对于墙角、立柱等环境结构的观测。SLAM算法每收到一帧有效的激光扫描，就会问自己：“如果我现在在这个位置，那么我看到的激光点应该和之前看到的哪些点重合？”通过把当前扫描与历史扫描或已有地图对齐，算法同时修正“我在哪”和“墙在哪”两个问题；当小车再次观察到去过的区域并形成可靠回环时，还能进一步修正长期累积误差。

相比只依靠里程计或IMU，激光SLAM能利用环境中的静态结构不断增加约束，因此通常可以减小累积漂移。但环境特征太少、雷达安装不准、时间不同步或扫描匹配错误，仍然会让地图出现问题。
\end{minipage}
}
\end{center}

在本章中，SLAM Toolbox就是做了这些事情——它读取\code{/scan}、局部里程计和TF树，输出\code{/map}并发布\code{map}到\code{odom\_combined}的变换。

\booksection{7.2 激光SLAM建图（以SLAM Toolbox为主线）}

上一节我们理解了SLAM的数学原理——把定位和建图变成一个优化问题，用位姿图和图优化来求解。但理论归理论，在实际的小车上我们到底怎么跑起来？

这一节，我们用\textbf{SLAM Toolbox}——ROS~2生态中常用的2D激光SLAM工具——在我们的麦轮小车上实际跑一遍建图。我们会一步步讲清楚：它是什么、怎么启动、参数怎么调、地图怎么保存。

\booksubsection{7.2.1 SLAM Toolbox是什么？它的输入和输出是什么？}

\begin{center}
\fbox{
\begin{minipage}{0.85\textwidth}
\small
\textbf{一句话先有个印象}

SLAM Toolbox就像一个“地图绘制员”——你推着小车在房间里走一圈，它一边听着雷达报数，一边用图优化把数据拼成一张完整的地图。
\end{minipage}
}
\end{center}

SLAM Toolbox是一个面向ROS和ROS~2的开源2D SLAM工具包，由Steve Macenski主导开发。它建立在SRI International开源的Open Karto库之上，并对测量匹配、图优化、序列化和运行模式进行了大量改进\cite{macenski2021slamtoolbox}。

\textbf{为什么选SLAM Toolbox而不是其他方案？}

市面上常见的2D激光SLAM方案有几种，各有特点：

\begin{table}[h]
\centering
\begin{tabular}{|c|c|c|}
\hline
\textbf{方案} & \textbf{核心特点} & \textbf{适用场景} \\
\hline
Gmapping & 基于粒子滤波，轻量 & 小场景、资源受限设备 \\
\hline
Hector SLAM & 纯激光匹配，不要求里程计，但没有回环闭合 & 高频激光、环境特征较丰富的场景 \\
\hline
Cartographer & 图优化、子地图和回环，系统较复杂 & 曾用于大场景实时建图；原谷歌项目已停止维护 \\
\hline
SLAM Toolbox & 图优化、多会话建图、会话序列化与交互工具 & \textbf{我们的选择}——ROS 2支持完善，便于教学 \\
\hline
\end{tabular}
\end{table}

对于我们的麦轮小车来说，选择SLAM Toolbox的理由很简单：\textbf{与ROS~2和Nav2集成良好，配置路径清楚，也便于观察和排查问题}。

SLAM Toolbox的内部工作流程可以分成四个层次，从下往上看：

\begin{enumerate}
\item \textbf{数据输入层}：接收激光雷达数据（\code{/scan}），并通过TF取得局部里程计变换和雷达外参。
\item \textbf{前端处理层}（Front-end）：对满足运动距离、转角和时间等更新条件的激光测量进行\textbf{扫描匹配}——把当前扫描和局部地图对齐，估计出相对位移W。
\item \textbf{图管理层}（Graph）：维护一个位姿图，节点是小车的各个历史位姿，边是里程计约束和扫描匹配约束。
\item \textbf{后端优化层}（Back-end）：当检测到回环或定期触发时，运行\textbf{图优化}，调整所有节点让全局误差最小。这一层跑得慢一些，但保证了地图的全局一致性。
\end{enumerate}

有了这个分层结构，SLAM Toolbox可以做到：前端实时跟踪（小车动到哪里，地图实时更新），后端定期整理（发现走歪了就拉回来）。

\textbf{输入：SLAM Toolbox需要什么？}

SLAM Toolbox的输入接口主要有3个核心部分：

\begin{enumerate}
\item \textbf{激光雷达数据}（通常为话题\code{/scan}）：提供环境的几何观测，是SLAM最主要的“眼睛”。
\item \textbf{局部里程计变换}：SLAM Toolbox通过TF查询\code{odom\_combined}$\rightarrow$\code{base\_footprint}，获得连续的局部运动估计。
\item \textbf{雷达外参}：TF树中必须存在\code{base\_footprint}$\rightarrow$\code{laser}，也就是告诉SLAM“激光雷达装在小车的哪个位置和方向”。
\end{enumerate}


\textbf{输出：SLAM Toolbox会给我们什么？}

SLAM Toolbox在运行时会发布两样关键的东西：

\begin{enumerate}
\item \textbf{地图}（话题：\code{/map}）：一张不断更新的栅格地图。随着小车探索，这张地图会越来越完整。
\item \textbf{全局位姿修正}（TF变换：\code{map→odom\_combined}）：告诉小车“你现在在地图上的什么位置”。这个变换每隔一段时间会更新一次，把累积的误差拉回来。
\end{enumerate}

在RViz中，你可以同时看到激光点云、正在构建的地图，以及小车的实时位置——非常直观。

\textbf{启动SLAM Toolbox建图}

在麦轮小车上启动建图模式，只需要一条命令：

\begin{verbatim}
# ROS 2 Humble；实物麦轮小车；来自wheeltec_robot_slam源码包
ros2 launch wheeltec_slam_toolbox online_sync.launch.py
\end{verbatim}

这个启动文件会做以下几件事：

\begin{enumerate}
\item 启动底盘驱动节点（发布\code{/odom\_combined}）；
\item 启动激光雷达驱动节点（发布\code{/scan}）；
\item 启动SLAM Toolbox的同步建图节点（\code{sync\_slam\_toolbox\_node}）；
\item 把默认的\code{/odom}重映射为\code{/odom\_combined}；
\item 加载配置文件\code{mapper\_params\_online\_sync.yaml}。
\end{enumerate}

启动之后，在RViz中打开\code{/map}话题显示，你应该能看到一张地图正在慢慢地被画出来。

\begin{center}
\fbox{
\begin{minipage}{0.85\textwidth}
\small
\textbf{第一次建图应该做什么}

第一次建图不要急着覆盖整个房间。你的目标应该是：\textbf{画出一个闭合回路}。

你需要找一个有墙角、有柱子、有家具——即有几何特征的小环境，推着小车慢慢走一圈。
因为只有形成了闭合回路，SLAM才能触发回环优化，修正累积误差。如果你只是来回走直线，或者走了一个“S”形而没有回到起点，那地图就会像一条歪歪扭扭的蛇——越远越偏，永远得不到一张平整的地图。
\end{minipage}
}
\end{center}

\begin{figure}[h]
    \centering
    \includegraphics[width=0.3\linewidth]{pictures/7.2地图.png}
    \caption{用SLAM Toolbox构建的地图}
    \label{fig:placeholder}
\end{figure}

\booksubsection{7.2.2 同步建图、异步建图与纯定位：有什么区别？怎么选？}

SLAM Toolbox提供三个主要运行模式：\textbf{同步建图}、\textbf{异步建图}和\textbf{纯定位}\cite{macenski2021slamtoolbox}。前两种模式回答“地图怎样建立”，第三种模式回答“已经有了SLAM会话，怎样只用它定位”。本章重点介绍前两种建图模式。

\textbf{同步建图（Sync Mode）}

同步模式按照顺序处理满足更新条件的激光测量，并保留等待加入SLAM问题的测量缓冲区。

\begin{itemize}
\item 如果处理速度跟得上：有效测量可以及时进入建图过程。
\item 如果复杂回环或图优化耗时较长：待处理测量可能产生延迟，因此需要足够的计算资源。
\end{itemize}

同步模式更重视测量和地图信息的完整性，也适合离线处理。对于我们的麦轮小车在室内环境建图，只要计算负载可控，就可以先把它作为首选。

\textbf{异步建图（Async Mode）}

异步模式则采取另一种策略：优先保证实时性。

\begin{itemize}
\item 如果处理速度跟得上：新的有效测量可以正常进入处理过程。
\item 如果上一项测量尚未处理完成：这段时间到达的部分有效测量可能不会进入地图；处理完成后，系统再接收满足更新条件的新测量。
\end{itemize}

异步模式的好处是不会因为一次复杂回环而持续落后于实时进度，小车移动时定位更新更容易跟上当前状态。代价是可能跳过部分扫描，地图信息不如同步模式完整。它更适合优先保证实时定位性能、不能接受处理过程持续滞后的场景。

\begin{center}
\fbox{
\begin{minipage}{0.85\textwidth}
\small
\textbf{一个帮助理解的比喻}

同步模式就像一个认真整理课堂笔记的学生——只要内容满足记录条件，他就按顺序整理；老师讲得太快时，等待整理的笔记会越来越多。

异步模式就像一个边听边拣重点记的学生——如果上一条还没整理完，期间的一些内容可能不会记入；整理完成后，再从新的有效内容继续。他的笔记可能不那么全，但不会一直落后于老师当前讲到的位置。
\end{minipage}
}
\end{center}

\textbf{如何切换到异步模式？}

把启动命令中的\code{online\_sync}换成\code{online\_async}即可：

\begin{verbatim}
ros2 launch wheeltec_slam_toolbox online_async.launch.py
\end{verbatim}

在配置文件方面，同步模式使用的是\code{mapper\_params\_online\_sync.yaml}，异步模式使用的是\code{mapper\_params\_online\_async.yaml}。



\booksubsection{7.2.3 关键参数详解与地图生成}

SLAM Toolbox的配置文件（\code{mapper\_params\_online\_sync.yaml}）里有好几十个参数。对于初学者来说，有两大类参数需要了解：一是“必配参数”，二是“可调参数”。

\textbf{必配参数（接口合同）}

这五个参数是SLAM正常运行的“准生证”，缺一不可：

\begin{verbatim}
odom_frame: odom_combined      # 里程计坐标系的名字
map_frame: map                 # 地图坐标系的名字
base_frame: base_footprint     # 机器人本体坐标系的名字
scan_topic: /scan              # 激光雷达数据的发布话题
mode: mapping                  # 运行模式：mapping（建图）或 localization（纯定位）
\end{verbatim}

这些参数在启动文件中已经配置好了。如果没有特殊原因（比如你给小车改了名字），不要去修改它们。

\textbf{可调参数（重点关注这几类）}

\subsubsection*{① 地图分辨率相关}

\begin{table}[h]
\centering
\begin{tabular}{|c|c|c|c|}
\hline
\textbf{参数名} & \textbf{含义} & \textbf{默认值} & \textbf{建议范围} \\
\hline
\code{resolution} & 每个像素代表多少米 & 0.05 & 0.02~0.10 \\
\hline
\end{tabular}
\end{table}

\code{resolution} 决定了地图的精细程度。0.05表示每个像素是5cm——也就是说，一面2m长的墙在地图上会用40个像素来表示。

\begin{itemize}
\item 值越小（如0.02）：地图越精细，能看到更小的物体细节，但计算量和内存占用会增大。
\item 值越大（如0.10）：地图越粗糙，小物体会被“抹掉”，但计算速度快、省内存。
\end{itemize}

对于大多数室内场景，0.05是很好的折中。如果你在仿真环境中跑，可以尝试0.03；如果你在资源有限的嵌入式平台上跑，可以用0.08。

(这里准备配一张不同值的对比图）
\subsubsection*{② 扫描匹配相关（影响运动跟踪质量）}

\begin{table}[h]
\centering
\begin{tabular}{|c|c|c|c|}
\hline
\textbf{参数名} & \textbf{含义} & \textbf{默认值} & \textbf{调整方向} \\
\hline
\code{minimum\_travel\_distance} & 移动多少米才添加新节点 & 0.5 & 越小→节点越密→地图越精细，但计算量越大 \\
\hline
\code{minimum\_travel\_heading} & 旋转多少弧度才添加新节点 & 0.5 & 同上 \\
\hline
\code{scan\_buffer\_size} & 缓存多少帧扫描用于匹配 & 10 & 越大→匹配更稳定，但内存更大 \\
\hline
\code{match\_search\_radius} & 匹配时的搜索半径 & 0.1 & 越大→能容忍更大的初始误差，但可能匹配错 \\
\hline
\end{tabular}
\end{table}

这里的“节点”指的是位姿图中的节点。节点越多，轨迹的描述就越精细，图优化的计算量也越大。\code{minimum\_travel\_distance} 这个参数控制着节点的密度：设得越小，节点越密，地图越精细；设得越大，节点越少，地图越粗糙但跑得更快。

对于我们的麦轮小车，默认值通常工作得很好。如果你发现地图中直线墙出现锯齿，可以试着把这两个参数调小一半。

\subsubsection*{③ 回环检测相关（影响地图的全局一致性）}

\begin{table}[h]
\centering
\begin{tabular}{|c|c|c|c|}
\hline
\textbf{参数名} & \textbf{含义} & \textbf{默认值} & \textbf{说明} \\
\hline
\code{do\_loop\_closing} & 是否启用回环检测 & true & 强烈建议保持开启 \\
\hline
\code{loop\_search\_maximum\_distance} & 搜索回环的最大距离 & 3.0 & 搜索范围太大→误检测；太小→漏检测 \\
\hline
\code{loop\_match\_minimum\_chain\_size} & 回环确认所需的最少匹配数 & 10 & 值越大→回环检测越保守（更少误报） \\
\hline
\end{tabular}
\end{table}

回环检测是SLAM能否完成精准建图的关键。如果不进行回环检测（\code{do\_loop\_closing}设为false），SLAM就会退化，只能修正局部误差，无法修正大尺度漂移。

\begin{center}
\fbox{
\begin{minipage}{0.85\textwidth}
\small
\textbf{回环检测是怎么工作的？}

SLAM Toolbox会从历史数据中搜索可能与当前位置重叠的候选区域，再对候选扫描执行匹配和约束检验。如果匹配结果满足回环条件，就在位姿图中增加一条回环边，然后触发图优化，把整条轨迹调整得更加一致。它不是简单地把当前位姿和所有历史位姿逐个硬比较。
\end{minipage}
}
\end{center}

\textbf{第一次跑图时，优先用默认参数}

新手常犯的一个错误是：第一次建图效果不理想，就开始疯狂调参。更好的做法是：

\begin{enumerate}
\item 先用默认参数跑一遍完整的建图流程；
\item 观察问题出在哪里——是地图整体扭曲？是局部锯齿？是回环没闭合？
\item 根据问题类型，有针对性地调整一两组参数；
\item 重新建图，比较效果。
\end{enumerate}
(这边想加入一些建图效果不好的图片，包含整体扭曲，局部锯齿和回环没闭合。）

记住：\textbf{参数只需要微调，数据质量才是真正的基础}。如果激光雷达装歪了、轮子打滑严重、小车走得太快，调再多的参数也没用。

\textbf{保存地图}

建图完成后，地图还存在SLAM Toolbox的内存里。我们需要把它保存成文件，供后续使用（比如下一章的导航）。

保存地图用这条命令：

\begin{verbatim}
ros2 run nav2_map_server map_saver -f ~/my_map
\end{verbatim}

这会生成两个文件：

\begin{enumerate}
\item \texttt{my\_map.pgm}：地图图片（PGM格式，灰色图片。黑色=障碍物，白色=空地，灰色=未知区域）。
\item \texttt{my\_map.yaml}：地图元数据文件（分辨率、原点位置、占用阈值等）。
\end{enumerate}

这里还要区分\textbf{栅格地图}和\textbf{序列化建图会话}。PGM和YAML供地图服务器、AMCL和导航使用，只保存占据栅格及其元数据，并不能恢复SLAM内部的历史扫描、约束和位姿图。如果以后还要继续扩展地图、进行多会话建图、执行SLAM Toolbox纯定位或手动调整位姿图，还应使用SLAM Toolbox的序列化功能保存本次建图会话。两类文件用途不同，不能互相替代。

保存好之后，你可以用图片查看器打开\code{.pgm}文件，看看地图画得怎么样——墙是否平直、角落是否清晰、有没有重影。

\begin{center}
\fbox{
\begin{minipage}{0.85\textwidth}
\small
\textbf{保存地图时的常见问题}

\textbf{问题：提示“No map received”}

原因：SLAM Toolbox还没有发布任何地图数据，或者网络通信有问题。解决方法：确认RViz中可以看到地图后，再运行保存命令。

\textbf{问题：保存的地图是空的或全黑的}

原因可能是地图尚未生成、保存工具没有正确收到\code{/map}、ROS域或命名空间不一致、建图节点已经退出，或者有效探索区域太小。解决时先用\code{ros2 topic echo /map --once}确认地图消息是否存在，再检查保存工具和建图节点是否处于同一通信环境，不要一上来就重新跑完整流程。
\end{minipage}
}
\end{center}

\textbf{SLAM Toolbox还能做什么？}

除了在线建图和纯定位，SLAM Toolbox还可以重新加载一次序列化会话，在下一次运行中继续完善或扩展位姿图，这就是\textbf{多会话建图}。它还提供位姿图交互调整和地图合并工具：遇到特别困难的回环时，用户可以在可视化界面中辅助对齐扫描；分别建立的序列化地图也可以合并为一张复合地图。

论文还介绍了原型性的终身建图和分布式建图应用。这里的“原型性”很重要：它说明这些方向已经可以演示和研究，但不能把它们直接理解成所有版本中都已成熟、默认开启的标准流程。

\textbf{本章的一个小练习}

在你成功建图并保存地图之后，试着回答这三个问题：

\begin{enumerate}
\item 你的地图中，墙的宽度大概占几个像素？（提示：看\code{resolution}参数）
\item 回到起点时，地图上有没有出现“双层墙”？（如果有，说明回环检测没生效）
\item 你的地图边缘是否平滑？（如果不平滑，可能是里程计噪声太大或扫描匹配参数需要调整）
\end{enumerate}

这三个问题的答案，会告诉你下一步应该调整哪些参数。



\booksection{7.3 建图质量优化与导航前检查}

地图能够保存、AMCL能够启动，并不等于建图与定位质量已经合格。进入第8章导航前，本节要建立一套更可靠的判断方法：先描述现象，再寻找原始证据，最后才调整算法参数。

\begin{center}
\fbox{
\begin{minipage}{0.85\textwidth}
\small
\textbf{质量优化不是把参数调到地图好看}

地图首先是定位和导航使用的数据，而不是一张展示图片。墙面平直但尺度错误，仍然不是合格地图；粒子云很集中但激光没有贴合墙面，仍然不是正确定位。每次优化都要回答三个问题：现象是什么？哪条数据能够证明原因？修改后哪项指标确实变好了？
\end{minipage}
}
\end{center}

建图与定位共享激光、里程计和TF，因此许多故障的外观相似。推荐始终按以下层次检查：

\begin{enumerate}
  \item \textbf{传感器层：}雷达距离是否稳定，时间戳是否连续，编码器与IMU方向是否正确；
  \item \textbf{坐标层：}雷达外参是否正确，TF主链是否连通且只有一个发布者；
  \item \textbf{运动层：}\code{odom\_combined}能否正确反映直行、横移和旋转；
  \item \textbf{算法层：}扫描匹配、回环检测、粒子滤波及相应参数是否合理；
  \item \textbf{任务层：}建图路线、初始位姿、场地动态变化是否为算法提供了足够信息。
\end{enumerate}

如果底层证据不可靠，直接调SLAM或AMCL参数只会暂时掩盖问题。

\booksubsection{7.3.1 地图重影、断裂与回环闭合}

同一堵墙在地图上出现两条平行边，称为重影。常见原因包括：里程计漂移过大、雷达外参错误、时间不同步、移动过快、环境特征不足或回环没有建立。

地图质量问题不能只用“好”或“不好”描述。观察地图时，至少区分下面四类现象。

\begin{table}[H]
\centering
\begin{tabular}{|p{0.19\textwidth}|p{0.31\textwidth}|p{0.38\textwidth}|}
\hline
\textbf{现象} & \textbf{地图表现} & \textbf{优先怀疑对象} \\
\hline
重影 & 同一堵墙形成两条或多条近似平行边 & 里程计累计漂移、雷达外参、时间不同步、回环未闭合 \\
\hline
断裂 & 本应连续的墙面中间缺失，或相邻区域没有连接 & 激光丢帧、有效量程设置、移动过快、扫描匹配失败 \\
\hline
弯曲与锯齿 & 真实直墙在地图中弯曲、呈波浪状或锯齿状 & 航向误差、振动、雷达安装松动、扫描畸变、角度分辨率 \\
\hline
整体错位 & 局部地图看起来正常，但闭环后房间无法重合 & 长距离累计误差、错误回环、环境重复特征、回环约束不足 \\
\hline
\end{tabular}
\end{table}

\textbf{重影不一定都是回环参数造成的}

先让小车保持静止，在RViz中观察连续多帧\code{/scan}。如果静止时激光端点本身就在抖动，问题可能来自雷达振动、数据噪声或时间戳；如果静止扫描稳定，而小车一转动墙面就变厚，应重点检查雷达相对\code{base\_footprint}的安装角度、旋转方向和时间同步；如果小范围地图正常，只有长距离返回起点时出现双层墙，再重点检查里程计累计漂移与回环检测。

雷达外参中的小角度误差会随测量距离放大。若雷达的实际偏航角与TF中记录的偏航角相差$\Delta\theta$，距离为$r$的激光点会产生近似横向偏差
\begin{equation}
e \approx r\sin(\Delta\theta).
\end{equation}
当$\Delta\theta$较小时，还可近似写成$e\approx r\Delta\theta$。这说明近处墙面看起来正常，并不能证明外参完全正确；在较长走廊中，小角度误差更容易表现为墙面错位。

\textbf{断裂要先区分“没有观测”还是“没有匹配”}

地图中一段墙消失，可能是雷达根本没有获得有效距离，例如玻璃透射、黑色吸光材料、超出量程或消息丢失；也可能是雷达看到了墙，但扫描匹配未能把该帧放到正确位置。前一种情况应检查\code{/scan}中的\code{range\_min}、\code{range\_max}、无穷值比例和发布频率，后一种情况则应对比同一时刻的里程计、TF与地图更新。

不要用“增大最大量程”来强行填补所有空白。雷达标称量程、室内反射条件和算法使用量程不是一回事。过远且不稳定的测量可能给扫描匹配引入错误约束。应先用静态障碍物在不同距离处验证有效测量范围，再决定配置值。

\textbf{回环闭合是全局一致性的检查点}

小车沿一条闭合路线回到起点附近时，SLAM Toolbox会尝试判断当前位置是否与历史位置对应。若匹配可信，位姿图中会增加回环约束，后端优化再把整段轨迹调整得更一致。回环的作用不是简单地把终点“拖回起点”，而是让沿途多个位姿共同满足里程计、扫描匹配和回环约束。

判断回环是否有效，可以观察：

\begin{itemize}
  \item 返回起点前后，同一墙面是否由双层逐渐合并；
  \item 地图是否发生一次整体但合理的微调，而不是局部突然撕裂；
  \item 回环后新采集的扫描能否继续稳定地贴合优化后的地图；
  \item 重复同一路线时，闭环结果是否具有一致性。
\end{itemize}

若算法把两个相似但不同的走廊误认为同一地点，就会产生错误回环。错误回环往往表现为地图被强行折叠、房间互相覆盖或原本平直的走廊突然弯折。此时盲目提高回环搜索范围会使问题更严重，应增加具有辨识度的观测、降低移动速度，并核对候选区域是否真的具有相同几何关系。

诊断时先看原始证据：静止时扫描是否稳定，直行时局部里程计是否合理，TF是否连续。只有底层正常后，才调整SLAM匹配参数。

推荐采用下面的小实验来定位地图问题：

\begin{enumerate}
  \item 原地静止\SI{10}{\second}，记录\code{/scan}稳定性；
  \item 低速直行一段已测量距离，比较实物距离与\code{odom\_combined}；
  \item 原地旋转一周，检查地图中的固定墙面是否增厚或旋转错位；
  \item 建立一个小矩形闭环，观察回到起点前后的地图变化；
  \item 小闭环可靠后，再扩大路线，不要直接在大场地中反复试错。
\end{enumerate}

每一步只增加一种运动或一种环境复杂度，才能判断问题从哪一步开始出现。

\booksubsection{7.3.2 导航前定位检查：不收敛与位姿跳变}

定位跳变可能来自两个完全不同的情况。一种是算法发现先前位置错了，进行合理的全局修正；另一种是错误匹配或传感器异常。区分二者要看修正后激光是否更贴合地图，以及跳变是否频繁反复。

定位质量至少包含三层含义：

\begin{enumerate}
  \item \textbf{正确性：}估计位置与机器人真实位置相符；
  \item \textbf{确定性：}粒子主要集中在一个合理区域，协方差没有持续增大；
  \item \textbf{稳定性：}小车正常运动时全局位姿变化合理，静止时不会反复跳到其他区域。
\end{enumerate}

粒子集中只说明算法“相信”某个答案，不保证答案正确。尤其在长直走廊、重复门洞、规则货架或外观相似的教室中，多个位置可能产生相似的激光轮廓。此时AMCL可能形成多个粒子簇，也可能错误地集中到其中一个位置。Navigation2的长时间实验也表明，重复特征、人员遮挡等因素会降低定位置信度；只要仍有足够激光束能够命中静态地图，定位通常更容易维持。

\textbf{合理修正与异常跳变的区别}

AMCL对\code{map}$\rightarrow$\code{odom\_combined}进行全局修正时，\code{map}坐标系中的机器人位置可能出现变化，而\code{odom\_combined}$\rightarrow$\code{base\_footprint}仍应保持局部连续。若修正后激光明显更贴合地图、粒子更加集中，并且之后能够稳定维持，这通常是合理修正。若位置在两个区域之间反复切换，或每次修正后激光反而偏离墙面，则更可能是错误匹配、重复环境、外参或时间问题。

\begin{table}[H]
\centering
\begin{tabular}{|p{0.20\textwidth}|p{0.31\textwidth}|p{0.37\textwidth}|}
\hline
\textbf{现象} & \textbf{可能原因} & \textbf{验证方法} \\
\hline
粒子长期分散 & 初始范围过大、环境特征少、激光与地图不一致 & 移到不对称墙角附近，重新设置初始位姿并缓慢旋转 \\
\hline
粒子集中但位置错误 & 初始位姿错误、场景重复、地图方向或尺度错误 & 同时检查实物位置、激光贴图和地图中的独特结构 \\
\hline
横移时明显跳动 & 使用差速模型、横向里程计错误、麦轮打滑 & 核对全向运动模型，单独执行低速横移实验 \\
\hline
人员经过后置信度下降 & 大量激光束被动态目标遮挡 & 人员离开后低速旋转，检查静态墙面观测是否恢复 \\
\hline
静止时持续漂移 & 里程计静态漂移、雷达数据抖动、TF时间异常 & 分别记录静止里程计、连续扫描和TF时间错误 \\
\hline
\end{tabular}
\end{table}

当人员或移动物体占据大量视野时，实际激光与建图时的静态地图会出现差异。AMCL并不要求每束激光都与地图完全一致，但若能够命中静态墙面的有效激光过少，错误位置与正确位置之间的权重差异就会减小。此时可以等待动态遮挡离开，或让小车在安全空间内低速旋转，以重新获得不同方向的静态环境特征。

\textbf{建立可重复的定位评价方法}

仅凭一次RViz画面不能比较参数。可以在地图中选取三个具有不同特征的位置：明显墙角、普通开阔区和重复走廊。每个位置重复若干次以下流程：放置小车、设置相同精度的初始位姿、执行相同低速动作、记录达到稳定粒子簇所需时间，并测量最终位置误差与航向误差。

若第$j$次实验的平面位置估计为$(\hat{x}_j,\hat{y}_j)$，参考位置为$(x_j,y_j)$，则位置误差可写为
\begin{equation}
e_{p,j}=\sqrt{(\hat{x}_j-x_j)^2+(\hat{y}_j-y_j)^2}.
\end{equation}
多次实验可计算平均位置误差
\begin{equation}
\bar{e}_p=\frac{1}{N}\sum_{j=1}^{N}e_{p,j}.
\end{equation}
参考位置可以来自场地测量标记、仿真真值或更高精度定位设备，但必须说明来源与精度。没有参考值时，可以报告重复定位的离散程度，但不能把重复性直接称为绝对精度。

\booksubsection{7.3.3 实物建图场地要求与运动路线建议}

实物建图建议遵循“慢、近、闭环”。速度保持较低，路线靠近有几何特征的区域，尽早形成小闭环，再逐步扩大地图。玻璃、长直走廊、拥挤人群和频繁移动的物体都会增加难度。

\textbf{建图前先检查场地}

理想的首次建图场地应具有稳定、静止且分布在多个方向的几何特征，例如墙角、柱子、门洞和固定柜体。以下场景需要特别处理：

\begin{itemize}
  \item \textbf{玻璃与镜面：}激光可能透射或产生异常反射，使地图出现缺口或虚假障碍；
  \item \textbf{长直走廊：}沿走廊方向的几何约束较弱，重复门洞还可能引起错误匹配；
  \item \textbf{大面积空旷区域：}有效回波较少，扫描之间缺乏足够重叠；
  \item \textbf{人群与移动物体：}动态观测不会长期对应静态地图，可能遮挡真正的墙面；
  \item \textbf{斜坡、门槛和振动地面：}二维雷达平面随车体俯仰或滚转变化，原本同一高度的墙面可能被扫到不同位置。
\end{itemize}


\textbf{路线设计遵循“小闭环扩展”}

第一次路线可以围绕一个房间或一个障碍物形成小矩形闭环。确认闭环后地图没有明显重影，再从闭环边缘向外探索，并尽量返回已经建好的区域。这样，新区域通过多个方向与旧区域连接，回环候选更容易获得共同视野。

路线规划可以分为四个阶段：

\begin{enumerate}
  \item \textbf{起点观察：}在起点静止数秒，让雷达获得稳定扫描；
  \item \textbf{局部闭环：}低速绕行一个小区域并返回起点附近；
  \item \textbf{逐步扩展：}每次只增加一块新区域，随后回到已建区域；
  \item \textbf{复查覆盖：}沿关键墙面再次低速通过，检查地图边界和连接处。
\end{enumerate}

在狭窄区域不要贴墙过近。雷达需要看到足够范围的环境，小车也需要保留转弯、横移和安全停车空间。经过门口时尽量保持平稳，不要在门框附近突然高速旋转，因为局部视野快速变化容易增加扫描匹配难度。

麦轮横移可能带来比前进更明显的滚子滑动和车体振动，所以建议先只用平稳的前进与转弯建立可靠地图，再用小幅横移测试全向底盘对定位和扫描匹配的影响。

\textbf{一次只改变一个变量}

建图效果受速度、更新阈值、分辨率、扫描缓存、回环参数和传感器质量共同影响。为了得到可解释结论，每轮实验只改变一组相关参数，其余条件尽量保持一致，并记录：

\begin{itemize}
  \item 地图名称、参数文件与提交版本；
  \item 场地、路线、地面、人员流动和照明等条件；
  \item 最大前向、横向和角速度；
  \item \code{/scan}与里程计频率、CPU负载和是否丢帧；
  \item 地图重影、闭环误差和保存结果；
  \item 仿真与实物的差异及尚未验证的假设。
\end{itemize}

记住，地图优化的目的是形成一张尺度合理、主要静态边界清楚、闭环一致，并能支持AMCL稳定定位和第8章路径规划的地图。

\booksection{7.4 本章小结与习题}

本章给第6章的局部运动加上了环境约束。SLAM Toolbox利用\code{/scan}以及TF中的局部里程计变换建立\code{/map}，地图可以保存为图像与YAML，SLAM会话还可以另行序列化。保存好的静态地图将作为第8章AMCL定位与自主导航的环境基准。

沿着本章主线，小车完成了三次能力升级。

\begin{enumerate}
  \item \textbf{从局部运动到全局环境：}第6章的里程计只能提供连续但会漂移的局部位姿；激光雷达把运动与固定环境联系起来。
  \item \textbf{从扫描到地图：}SLAM Toolbox通过扫描匹配、位姿图和回环优化，一边估计轨迹，一边建立占据栅格地图。
  \item \textbf{从内存地图到可复用文件：}地图图像保存栅格，YAML保存分辨率、原点和阈值等元数据，两者共同描述导航地图；SLAM Toolbox序列化会话则保存历史扫描、约束和位姿图，用于继续建图、多会话建图或纯定位。
\end{enumerate}

本章最重要的系统关系不是某一个参数，而是
\begin{equation}
\code{map}\rightarrow\code{odom\_combined}\rightarrow\code{base\_footprint}\rightarrow\code{laser}.
\end{equation}
建图时，SLAM Toolbox负责\code{map}$\rightarrow$\code{odom\_combined}；第6章的局部状态估计负责\code{odom\_combined}$\rightarrow$\code{base\_footprint}，机器人模型或静态TF负责\code{base\_footprint}$\rightarrow$雷达坐标系。进入第8章的静态地图导航后，这条全局变换改由AMCL负责，同一时刻不能让两者重复发布。

\begin{center}
\fbox{
\begin{minipage}{0.85\textwidth}
\small
\textbf{进入第8章前的建图交付条件}

小车应能建立包含闭环的室内地图并正确保存；PGM与YAML应成对存在，分辨率、原点和路径可核查；重新启动后，地图服务器能够加载同一地图；直行、横移和旋转时，局部里程计方向正确且TF保持连续。初始位姿设置与AMCL收敛不再作为“进入导航章之前已经完成”的前提，而是第8章要明确执行和验证的第一段导航流程。
\end{minipage}
}
\end{center}

\booksubsection{7.4.1 SLAM与坐标系统分析题}
\begin{enumerate}
  \item 用自己的话解释“同时定位与建图”中的相互依赖关系。为什么未知地图中不能只先完成定位，再开始建图？
  \item 解释\code{map}与\code{odom\_combined}为什么不能简单合并成一个坐标系。分别说明哪个坐标系允许全局修正，哪个坐标系应保持局部连续。
  \item 画出\code{map}$\rightarrow$\code{odom\_combined}$\rightarrow$\code{base\_footprint}$\rightarrow$\code{laser}的TF链，并标出SLAM、AMCL、EKF和静态TF各自负责的变换。
  \item 说明PGM与YAML分别保存什么。如果只复制PGM、不复制YAML，地图尺度和原点信息会发生什么问题？
  \item 已知地图分辨率为\SI{0.05}{\metre}，一堵墙在图像中长80个像素，估算其实际长度。若原点为$[-2.0,-1.0,0]$，解释前三个数分别表示什么。
  \item 比较SLAM Toolbox建图模式与AMCL定位模式的输入、输出和地图是否更新。为什么两个节点不应同时发布同一条\code{map}$\rightarrow$\code{odom\_combined}？
  \item 按“撒豆子”的比喻，分别说明粒子的运动更新、观测加权与重采样对应什么操作。粒子集中是否一定代表定位正确？为什么？
  \item 麦轮小车横移时，AMCL若误用差速运动模型，可能出现哪些现象？应通过哪三个最小运动实验验证？
\end{enumerate}

\booksubsection{7.4.2 建图、地图保存与导航定位衔接任务}

本任务分为仿真基线和实物复现两个阶段。不得直接使用他人已经完成的地图代替建图过程。

\textbf{阶段一：仿真建图基线}

\begin{enumerate}
  \item 启动麦轮小车、激光雷达、局部状态估计和SLAM Toolbox，记录节点、话题与TF主链；
  \item 检查\code{/scan}和\code{/odom\_combined}频率，分别执行低速直行、横移和旋转；
  \item 设计一条包含明显墙角的小闭环路线，先完成局部闭环，再逐步扩大探索范围；
  \item 在回到起点前后分别保存RViz观察结果，说明地图是否发生全局调整；
  \item 保存地图图像与YAML，并记录地图分辨率、原点和文件路径。
\end{enumerate}

\textbf{阶段二：重新启动AMCL定位}

\begin{enumerate}
  \item 完全停止在线建图节点，使用地图服务器加载刚保存的地图；
  \item 启动AMCL，确认麦轮小车选择全向运动模型，并检查\code{/map}、\code{/scan}、\code{/amcl\_pose}和\code{/particle\_cloud}；
  \item 在RViz中设置近似初始位姿，让小车缓慢旋转或短距离移动；
  \item 记录粒子由分散到集中的过程、激光与地图的重合情况，以及\code{map}$\rightarrow$\code{odom\_combined}；
  \item 让小车依次直行、横移和旋转后回到起点附近，检查全局位姿方向和稳定性。
\end{enumerate}

\textbf{阶段三：实物复现与差异分析}

在实物小车上重复相同任务。实物阶段优先降低速度与加速度，并记录地面、轮胎状态、雷达安装、人员流动和计算平台负载。对比仿真与实物时，至少回答：

\begin{itemize}
  \item 哪些话题、消息类型和TF接口保持不变？
  \item 哪些速度、运动噪声、激光范围或更新阈值需要调整？
  \item 实物地图出现了哪些仿真中没有的噪声、重影或断裂？
  \item 横移对里程计、扫描匹配和AMCL粒子分布有什么影响？
  \item 哪些结论已经运行验证，哪些仍然只是静态审查或推测？
\end{itemize}

最低提交物包括：启动命令与参数文件、接口清单、TF主链、建图路线、地图图像和YAML、回环前后证据、初始位姿、定位收敛证据、一次失败记录，以及仿真与实物差异表。只提交最终地图截图，不能证明完整流程已经正确运行。

\booksubsection{7.4.3 建图与定位故障诊断题}

下面各题要求先写“观察证据”，再写“可能原因”和“处理措施”。不能只回答“调整参数”或“重新启动”。

\begin{enumerate}
  \item \textbf{双层墙：}小车回到起点附近后，同一堵墙出现两条平行边。请从静止激光、雷达外参、旋转实验、里程计累计误差和回环检测五个方面安排检查顺序。
  \item \textbf{地图断裂：}小车经过玻璃门后，地图墙面出现缺口。怎样判断是激光无有效回波，还是扫描匹配把该帧放错了位置？应查看哪些原始字段和同步数据？
  \item \textbf{错误回环：}两个外观相似的走廊被地图错误地叠在一起。为什么简单增大回环搜索距离可能使问题恶化？如何通过路线和环境特征降低歧义？
  \item \textbf{粒子不集中：}AMCL重新启动后，粒子长期分散。请从地图、\code{/scan}、TF、初始位姿、动态遮挡和粒子参数六个方面列出检查顺序，并说明为什么粒子数量应最后调整。
  \item \textbf{错误收敛：}粒子已经集中，但机器人模型显示在隔壁房间。怎样证明这是错误收敛，而不是RViz显示问题？如何选择更有辨识度的动作重新定位？
  \item \textbf{横移跳变：}前进和旋转定位正常，只有横移时\code{/amcl\_pose}明显跳动。请检查AMCL运动模型、$v_y$里程计、麦轮打滑和激光时间戳，并设计一个最小复现实验。
  \item \textbf{静止漂移：}小车完全静止，但地图坐标中的位姿仍缓慢变化。分别说明如何判断误差来自局部里程计、激光抖动还是AMCL错误匹配。
\end{enumerate}

\textbf{综合故障注入实验}

从以下故障中选择一种进行可恢复的仿真实验：把雷达偏航外参加入一个小误差、把横向里程计符号反转、让部分节点使用错误时钟，或给AMCL选择差速运动模型。记录故障注入前后的\code{/scan}、\code{/odom\_combined}、TF、地图或粒子分布，按照本节的分层顺序找到原因，再恢复正确配置。报告中必须写明故障是人为注入，不得把未经运行的预期现象写成实测结果。

\paragraph{Technical English}
阅读并解释下面三句话：

\begin{quote}
``Loop closure improves the global consistency of a pose graph.''

``A localization system corrects odometric drift by matching sensor observations against a known map.''

``A concentrated particle cloud does not necessarily imply a correct pose estimate.''
\end{quote}

请指出\emph{loop closure}、\emph{global consistency}、\emph{odometric drift}、\emph{known map}和\emph{particle cloud}在本章中分别对应哪个概念或可观察现象，并用英文写一段不少于80词的建图或定位故障报告。报告至少包含现象、证据、原因假设和下一步检查。

% ============================================================
\bookchapter{第8章 自主导航与任务执行}

\section*{本章任务：给出目标点，让小车自己到达}

现在，小车已经有了地图，但重新启动后并不会自动知道自己在地图中的位置。自主导航的第一步不是发送目标，而是加载地图、给出带不确定性的初始位姿，并让AMCL利用激光与里程计完成定位收敛。只有确认机器人在\code{map}坐标系中的位姿可信后，才能安全地规划路径、跟踪轨迹并在失败时执行恢复。

Nav2是一组协作服务器。行为树决定下一步做什么，规划器计算全局路径，控制器把路径变成速度，代价地图表示哪里危险，行为服务器在受阻时执行等待、后退或清图等恢复动作。这样的设计把整体的导航任务流程和每个部分的具体算法分开，这样更换规划器不必重写整套任务逻辑，增加一种恢复动作也不必改动底盘驱动\cite{macenski2020marathon}。

本章仍然围绕同一辆麦轮小车展开。我们先在仿真中建立一条可重复的成功基线，再把同一组接口迁移到实物。完成本章后，读者应能够：

\begin{itemize}
  \item 说明地图加载、初始位姿设置和AMCL收敛为何是导航启动流程的必要步骤；
  \item 画出从目标位姿到\code{/cmd\_vel}的Nav2数据流，并说明各服务器的职责；
  \item 区分路径（path）与轨迹（trajectory），说明二者分别回答“经过哪里”和“怎样随时间运动”；
  \item 区分全局代价地图和局部代价地图，解释静态层、障碍层与膨胀层；
  \item 根据麦轮底盘的全向运动能力选择规划器和控制器，并配置速度、加速度及机器人轮廓；
  \item 读懂一个包含规划、控制和恢复分支的行为树；
  \item 用可观察证据定位“没有路径”“有路径但不走”和“遇障碍后恢复失败”等问题；
  \item 用成功率、路径长度、任务时间和恢复次数评价导航效果。
\end{itemize}

\noindent\statusstatic 本章的Nav2架构与算法说明主要依据Navigation2系统论文和ROS~2移动机器人算法综述\cite{macenski2020marathon,macenski2023survey}；项目包名、启动文件和参数来自本地源码。论文反映其发表时的设计与算法状态，具体到ROS~2 Humble的插件名称、参数和默认值，仍须以本地安装版本及实机实验为准。

\booksection{8.1 导航的起点：基于地图的全局定位（AMCL）}

上一节完成了两件事：先让小车在环境中行驶并建立地图，再把地图保存为图像和YAML文件。地图一旦固定，问题就从“边走边画”变成了“拿着一张已经画好的图，判断自己站在哪里”。这就是基于地图的全局定位。

本节使用Nav2中的AMCL完成这一任务。AMCL把地图、激光和里程计放在一起比较，估计小车在\code{map}坐标系中的位姿，并用\code{map}到\code{odom\_combined}的变换修正长期累计误差。


\booksubsection{8.1.1 粒子滤波定位思想}

建图阶段，地图和位置都未知。定位阶段，地图已经保存，只需要估计小车在地图中的位姿。

自适应蒙特卡洛定位（Adaptive Monte Carlo Localization，AMCL）使用一群“位置假设”，通常称为粒子。每个粒子都猜测小车可能在哪里。算法比较该位置下预期看到的激光与实际\code{/scan}，让解释得更好的粒子获得更高权重。

假设平面中的第$i$个粒子写成
\begin{equation}
\bm{x}^{[i]}_t=
\begin{bmatrix}
x^{[i]}_t & y^{[i]}_t & \theta^{[i]}_t
\end{bmatrix}^{\mathrm T},
\end{equation}
并带有权重$w^{[i]}_t$。粒子的坐标表示一种位姿猜测，权重表示该猜测与当前观测相符的程度。全部粒子共同近似机器人位姿的概率分布
\begin{equation}
p(\bm{x}_t\mid \bm{z}_{1:t},\bm{u}_{1:t},M),
\end{equation}
其中$M$是已知地图，$\bm{u}_{1:t}$是截至当前的里程计运动信息，$\bm{z}_{1:t}$是激光观测。

粒子滤波的一次更新可以分成三步。

\begin{enumerate}
  \item \textbf{运动更新：}根据两次里程计之间的位移，把每个粒子向前传播。为了表示打滑、编码器误差和模型误差，传播时会加入随机扰动。因此即使粒子原来集中在一点，移动后也会散开一些。
  \item \textbf{观测更新：}把实际激光束变换到每个粒子所猜测的位姿上，再与占据栅格地图比较。若激光端点大多落在地图障碍附近，该粒子的权重就较高；若激光穿过墙壁或落在大片空闲区域，权重就较低。
  \item \textbf{重采样：}按权重重新抽取粒子。高权重粒子更可能被保留和复制，低权重粒子逐渐消失。多轮更新后，粒子通常会在正确位姿附近形成一个紧凑粒子簇。
\end{enumerate}

观测更新可用一个简化公式表示：
\begin{equation}
w^{[i]}_t \propto
p\!\left(\bm{z}_t\mid \bm{x}^{[i]}_t,M\right).
\end{equation}
这个公式表示的意思是：“如果机器人真的位于这个粒子所表示的位置，当前激光出现的可能性有多大”。这是地图提供全局约束的方式。

\begin{center}
\fbox{
\begin{minipage}{0.85\textwidth}
\small
\textbf{用撒豆子类比粒子滤波}

可以把粒子滤波想成在地图上撒一把豆子。每颗豆子代表小车可能处于的一个位置和朝向。开始时，由于我们还不确定小车在哪里，豆子会散落在一片较大的区域。小车移动后，根据里程计记录的运动，把每颗豆子向相应方向拨动一点；考虑到轮子打滑和测量误差，各颗豆子的移动量还会略有不同。

接着，把每颗豆子所在位置“应该看到的墙”与雷达真正看到的墙进行比较。越符合实际观测的豆子越可信，可以保留并在它附近重新撒下更多豆子；明显对不上地图的豆子则被逐渐淘汰。经过多轮“移动、比较、筛选和重新撒豆子”，豆子会从大范围分散逐渐集中到正确位置附近，这就是定位收敛。

\end{minipage}
}
\end{center}

AMCL中的“自适应”主要体现在粒子数量可以随位姿分布的复杂程度变化。当位置已经比较确定时，不必始终保留大量粒子；当假设分散或环境容易混淆时，则需要更多粒子描述多个可能区域。这样可以在计算量与分布表达能力之间取得折中。需要注意，自适应并不意味着粒子越多越准确。传感器外参、时间戳或地图本身有错时，增加粒子只会让算法用更多计算量解释错误数据。


\booksubsection{8.1.2 AMCL的输入：地图、激光与里程计}

AMCL需要地图、激光扫描、局部里程计和完整TF。这四项不是可以互相替代的“参考信息”，而是一条缺一不可的数据链。

\begin{enumerate}
  \item \textbf{静态地图：}地图服务器读取上一节保存的YAML和PGM文件，并发布\code{/map}。地图给出每个栅格的占用状态、分辨率和原点。AMCL只在这张已知地图上定位，不负责继续填充未知区域。
  \item \textbf{激光扫描：}\code{/scan}提供各个角度的距离观测。AMCL通过\code{base\_footprint}到雷达坐标系的外参，把激光束放到每个粒子对应的地图位置进行匹配。
  \item \textbf{局部里程计：}AMCL通常不是简单订阅一个里程计话题来获得运动，而是通过TF查询\code{odom\_combined}到\code{base\_footprint}的相对变化。它关心两次观测之间小车怎样移动，而不是把局部里程计当成永远正确的全局坐标。
  \item \textbf{时间一致的TF：}至少要能连通\code{map}$\rightarrow$\code{odom\_combined}$\rightarrow$\code{base\_footprint}$\rightarrow$雷达坐标系。激光消息的时间戳必须落在TF缓存能够查询的范围内。
\end{enumerate}

AMCL估计机器人在\code{map}中的位姿后，会结合已有的\code{odom\_combined}$\rightarrow$\code{base\_footprint}计算并发布\code{map}$\rightarrow$\code{odom\_combined}。这样既保留局部里程计的连续性，又让全局位置能够被地图修正。不能让SLAM Toolbox和AMCL同时发布同一条\code{map}$\rightarrow$\code{odom\_combined}变换；进入定位模式前，应停止在线建图节点或关闭其中的TF发布职责。

AMCL常见的可观察输出包括：

\begin{itemize}
  \item \code{/amcl\_pose}：带协方差的全局位姿估计；
  \item \code{/particle\_cloud}：用于观察粒子在地图中的分布；
  \item \code{map}$\rightarrow$\code{odom\_combined}：提供给Nav2及其他全局任务使用的坐标变换。
\end{itemize}

在本案例中，启动定位前先确认地图文件存在，再由地图服务器与AMCL加载。具体启动文件名称取决于教学工程，下面给出接口层面的检查方法：

\begin{verbatim}
# ROS 2 Humble；命令用于接口检查；待教学平台运行验证
ros2 topic echo /map --once
ros2 topic hz /scan
ros2 topic echo /amcl_pose
ros2 run tf2_ros tf2_echo odom_combined base_footprint
ros2 run tf2_ros tf2_echo base_footprint laser
\end{verbatim}

如果实物雷达坐标系不叫\code{laser}，最后一条命令应替换为实际的\code{frame\_id}。可以先用\code{ros2 topic echo /scan --once}查看激光消息头，再核对TF。

\textbf{麦轮小车必须使用全向运动模型}

普通差速小车主要前进和转弯，理想情况下没有独立横向速度；麦轮小车却可以直接横移。因此，AMCL的机器人运动模型应选择全向模型，而不是差速模型。ROS 2 Humble的Nav2 AMCL配置通常写成：

\begin{verbatim}
amcl:
  ros__parameters:
    base_frame_id: base_footprint
    odom_frame_id: odom_combined
    global_frame_id: map
    scan_topic: scan
    robot_model_type: nav2_amcl::OmniMotionModel
\end{verbatim}

若错误使用差速模型，前进和原地旋转可能看起来正常，但横移时预测模型无法正确解释运动，粒子会不必要地扩散，激光匹配与里程计预测相互“打架”，最终表现为位姿滞后、跳动或局部错配。

全向模型的运动噪声通常由\code{alpha1}至\code{alpha5}等参数描述，它们反映旋转对旋转、平移对旋转、平移对平移以及横向平移等误差传播。这里不建议先背参数含义，更不建议一次修改全部参数。先通过直行、横移、旋转三个低速实验确认里程计方向与量级，再根据实际重复性调整运动噪声。

\textbf{激光模型如何给粒子打分？}

AMCL常见的激光观测模型包括波束模型（Beam Model）和似然场模型（Likelihood Field Model）。波束模型从粒子位置向地图中进行射线计算，再把预测距离与实测距离比较；似然场模型预先计算每个栅格到最近障碍物的距离，激光端点越靠近地图障碍，得分通常越高。室内栅格地图常使用似然场模型，因为它对小范围位姿误差更平滑，计算也较方便。

配置观测模型时，\code{laser\_model\_type}决定模型类型，\code{max\_beams}控制每帧用于计算的激光束数量，\code{sigma\_hit}、\code{z\_hit}和\code{z\_rand}等参数描述正确命中与随机测量的相对影响。束数过少会丢失环境约束，束数过多则增加计算量；权重设置不合理会让算法过度相信或忽略雷达。第一次运行应从工程默认值开始，而不是凭地图效果随意放大某一个权重。

\textbf{粒子数量与更新条件}

\code{min\_particles}和\code{max\_particles}给出自适应粒子数的上下限，\code{pf\_err}和\code{pf\_z}影响自适应采样的误差与置信要求，\code{update\_min\_d}和\code{update\_min\_a}决定移动或转动多少后进行一次滤波更新。较小的更新阈值能够更快利用新观测，但会增加计算负担，也可能反复处理变化很小的扫描。

对本案例而言，调参顺序应是：先保证地图、\code{/scan}、外参、局部里程计和时间戳正确；再确认全向运动模型；最后才观察CPU占用、粒子分布和定位误差，调整粒子数、更新阈值及激光模型参数。

如果AMCL不收敛，按下面顺序检查：地图是否加载、\code{/scan}是否有效、TF链是否完整、初始位姿是否合理。不要一开始就增加粒子数。

\begin{table}[H]
\centering
\begin{tabular}{|p{0.22\textwidth}|p{0.31\textwidth}|p{0.35\textwidth}|}
\hline
\textbf{检查对象} & \textbf{应观察的证据} & \textbf{常见异常} \\
\hline
地图 & \code{/map}有数据，分辨率和原点与保存时一致 & 地图未加载、路径错误、使用了另一张地图 \\
\hline
激光 & 频率稳定，距离有效，激光点能与墙面重合 & 全零或无穷值过多、雷达倒置、外参方向错误 \\
\hline
局部里程计 & 直行、横移和旋转方向正确，短时连续 & 横移符号错误、旋转尺度不对、静止时明显漂移 \\
\hline
TF与时间 & 主链唯一、连续，激光时间可查询到变换 & 多节点发布同一TF、外推到过去或未来、坐标系断开 \\
\hline
运动模型 & 麦轮底盘使用全向模型 & 横移时粒子异常扩散或位姿跳动 \\
\hline
\end{tabular}
\end{table}

\booksubsection{8.1.3 初始位姿设置与定位收敛判定}

AMCL刚启动时并不知道小车在地图哪里。所以为了防止计算量过大，我们通常会根据对现实中小车位置的观测，手动给它设定一个大致的初始位姿。设置初始位姿相当于告诉它：“大约在这里，朝向大约是这样。”

初始位姿不只包含$x$、$y$和航向角，还包含对这些量不确定程度的描述。在RViz中使用“2D Pose Estimate”工具时，鼠标按下的位置给出$x$、$y$，拖动箭头的方向给出航向。若只选对位置但箭头方向相反，机器人可能在几何重复的环境中长期匹配到错误墙面。

也可以向\code{/initialpose}发布\code{PoseWithCovarianceStamped}消息。无论使用哪种方式，都要保证消息的参考坐标系是\code{map}，时间戳与系统时钟一致。协方差太小相当于过度确信人工标注，真实位置若偏离较大，正确粒子可能一开始就很少；协方差太大则会让粒子铺得过散，收敛速度下降。课堂实验中先在RViz中给出合理范围，再观察粒子云，而不是直接把协方差设为零。

\begin{center}
\fbox{
\begin{minipage}{0.85\textwidth}
\small
\textbf{初始位姿不是“设置真值”}

在地图上点一下，只是给AMCL一个带不确定性的先验范围，并没有强行确定机器人就在该处。之后的激光观测仍会重新给粒子加权。如果初始箭头与真实朝向差别不大，粒子会逐渐收紧；如果初始位置明显错误，算法可能收敛到错误位置，也可能一直保持多个粒子簇。
\end{minipage}
}
\end{center}

设置后，缓慢旋转或移动小车，让雷达看到多个方向的环境特征。若粒子逐渐集中、激光点与地图墙面重合，说明定位正在收敛。若粒子集中在错误位置，可能是初始位姿、地图方向或激光外参错误。

仅凭粒子数量变少还不能宣布定位成功。至少应从四类证据判断收敛。

\begin{enumerate}
  \item \textbf{粒子分布：}\code{/particle\_cloud}由大范围分散逐渐变为正确位置附近的单个紧凑粒子簇。若形成两个相距很远的粒子簇，说明环境仍有歧义。
  \item \textbf{激光与地图重合：}在RViz中把固定坐标系设为\code{map}，同时显示地图、LaserScan和机器人模型。静止时，大部分打到静态墙面的激光端点应贴近地图障碍边界。
  \item \textbf{位姿与协方差：}\code{/amcl\_pose}的位置和航向应趋于稳定，其协方差总体减小并保持在可接受范围。协方差没有一个适用于所有平台的固定合格值，应通过同一场地的重复实验建立基线。
  \item \textbf{运动一致性：}让小车低速前进、横移、旋转，再回到起点附近。全局轨迹应与实际运动方向一致，停车后不应持续漂移或在多个相似位置之间反复跳变。
\end{enumerate}

可以把首次定位实验分成四步进行：

\begin{enumerate}
  \item 把小车放在地图中容易辨认的位置，例如不对称墙角附近，确认实物朝向；
  \item 在RViz中设置初始位姿，等待几秒时间，观察粒子是否覆盖真实位置；
  \item 以低角速度原地旋转一周或进行小范围“旋转—直行—横移”运动，让雷达观察多个方向；
  \item 保存收敛前后粒子分布、\code{/amcl\_pose}和激光贴图证据，并记录所用地图与参数文件。
\end{enumerate}

这里强调“缓慢”是因为高速旋转容易同时放大运动模型误差、雷达扫描畸变和TF时间误差。若慢速能够收敛而快速不能，先检查数据频率和时间链，不要立即把问题归因于粒子数量。

\textbf{机器人被搬动后怎么办？}

小车被抬起并放到地图另一处，而局部里程计没有感知这段运动，这种现象常称为“绑架机器人问题”。此时原有粒子簇仍会停留在旧位置，实际激光与地图突然不匹配。遇到这种情况时，可以重新设置初始位姿，或在安全条件下触发全局重定位，让粒子在更大地图范围内重新分布。但是全局初始化的搜索范围大、计算量高，而且对长直走廊等重复环境仍可能存在歧义，因此教学实验优先采用人工给出近似初始位姿的方法。

\textbf{集中在错误位置不等于正确收敛}

在规则教室、仓库货架或长直走廊中，不同区域可能具有相似几何结构。粒子会在一个错误但“看起来也能解释激光”的位置集中，这叫错误收敛，会导致定位直接失败。所以判断定位是否成功时必须同时查看粒子、激光贴图和实物位置。增加一次能看到不对称墙角、门洞或柱子的运动，通常比原地反复旋转更有信息量。

若定位失败，建议按以下顺序处理：

\begin{enumerate}
  \item 确认加载的是当前场地的正确地图，且地图方向、尺度和原点没有改变；
  \item 静止检查激光与机器人模型的相对位置，排除雷达外参错误；
  \item 分别做直行、横移和旋转，确认\code{odom\_combined}运动方向正确；
  \item 检查TF主链与时间戳，确认没有重复发布者和外推错误；
  \item 在有明显几何特征的位置重新设置初始位姿并慢速运动；
  \item 只有前五步都有证据支持时，才调整粒子数、运动噪声和激光模型参数。
\end{enumerate}

完成这一节后，系统应形成清晰的职责分工：地图服务器提供不变的环境模型，EKF提供连续局部运动，AMCL用激光把局部运动锚定到地图，TF把三者连接起来。本章后续的全局规划使用\code{map}坐标系，局部控制仍依赖连续的\code{odom\_combined}坐标系。正是这种“全局可修正、局部要连续”的分层设计，让导航既知道长期位置，又不会因一次定位修正而突然改变底盘控制速度。

\booksection{8.2 Nav2导航系统概述}

\booksubsection{8.2.1 自主导航数据流与计算图}

第一次接触Nav2时，很容易把它想成一个很大的程序：输入目标，程序在里面算一算，然后小车就开始走。真正运行起来却不是这样。Nav2更像一个分工明确的小组，每个人只负责一类工作，再由行为树安排他们什么时候出场。

可以把一次导航想成四个人合作。

\begin{itemize}
  \item \textbf{行为树导航器 (Behavior Tree Navigator / BT Navigator)：}接收目标，安排“规划、跟随、检查、恢复”的顺序，它本身不执行具体的规划或控制逻辑；
  \item \textbf{规划器服务器(Planner Server)：}它负责全局路径规划。它接收起点和终点，在全局代价地图 (Global Costmap) 上计算出一条从起点到终点的、尽可能优化的全局路径；
  \item \textbf{控制器服务器(Controller Server)：}它负责局部路径跟踪与控制。它接收规划器产生的全局路径，并结合局部代价地图 (Local Costmap) 上的实时障碍物信息，计算出机器人每一时刻应当执行的线速度和角速度(即\code{/cmd\_vel}信息）；
  \item \textbf{行为服务器(Behavior Server)：}它负责执行各种行为，尤其是恢复行为 (Recovery Behaviors)当路径被堵住或控制失败时，执行恢复动作。
\end{itemize}

\begin{center}
\fbox{
\begin{minipage}{0.85\textwidth}
\small
\textbf{先记住一句话}

规划器像看整张地图安排路线的人，控制器像握着方向盘观察车前路况的人，行为树则像任务负责人。负责人不会亲自给出全局路径，也不会改变电机转速；它只决定现在该让谁工作、失败后该找谁处理。
\end{minipage}
}
\end{center}

把定位与任务执行连起来，一次完整导航的大致数据流是：

\begin{center}
加载静态地图 $\rightarrow$ 设置初始位姿 $\rightarrow$ AMCL收敛 $\rightarrow$ 获得当前全局位姿\\
$\downarrow$\\
发送目标位姿 $\rightarrow$ 行为树导航器 $\rightarrow$ 全局规划器 $\rightarrow$ 全局路径\\
$\downarrow$\\
定位与TF $+$ 局部代价地图 $+$ 全局路径 $\rightarrow$ 控制器 $\rightarrow$ \code{/cmd\_vel} $\rightarrow$ 底盘
\end{center}

这条链里还有两路持续更新的信息。AMCL提供\code{map}$\rightarrow$\code{odom\_combined}，第6章的状态估计提供\code{odom\_combined}$\rightarrow$\code{base\_footprint}；激光雷达则通过\code{/scan}更新障碍信息。任何一条TF断开、时间戳过期或坐标系写错，都可能让后面的规划和控制停止。

\booksubsection{8.2.2 Nav2核心节点与生命周期管理}

Nav2由多个功能包和可替换插件组成。\code{nav2\_planner}提供规划服务器，\code{nav2\_controller}提供控制服务器，\code{nav2\_costmap\_2d}维护代价地图，\code{nav2\_bt\_navigator}执行行为树，\code{nav2\_amcl}完成基于地图的定位，\code{nav2\_map\_server}加载地图。

这里的“服务器”指的是职责清楚的ROS~2节点。规划服务器可以在运行时加载NavFn、Smac等规划插件，控制服务器可以加载DWB、RPP或其他控制插件。行为树中的动作节点通过Action调用这些服务器，服务器再把任务交给选定插件。

这些节点启动以后并不会立刻工作。因为如果所有节点一启动就抢着接收任务，错误会同时从多处冒出来，很难判断谁先出了问题。因此，Nav2大量使用生命周期节点。一个节点通常经历未配置、未激活、已激活和结束等状态：

\begin{enumerate}
  \item \textbf{配置：}读取参数、创建插件和通信接口，但还不执行导航；
  \item \textbf{激活：}开始接收目标、更新代价地图或输出控制；
  \item \textbf{停用与清理：}停止任务并释放相应资源；
  \item \textbf{错误处理：}某个组件异常时，由生命周期管理器协调停用或恢复。
\end{enumerate}

所以，节点进程存在不代表导航已经可用。如果RViz能看到节点名却无法发送目标，应继续检查生命周期状态、地图是否加载、定位是否激活以及各服务器是否都完成配置。

实例工程包含\code{navigation2-humble}源码快照，同时通过\code{wheeltec\_robot\_nav2}提供面向具体车型的启动与参数文件。阅读项目时要区分“Nav2通用实现”和“mini\_mec车型配置”。

\booksubsection{8.2.3 NavigateToPose Action任务接口}

现在所有服务器都准备好了，接下来怎样告诉小车去教室门口呢？单目标导航使用\code{NavigateToPose}动作（Action）给小车发布目标点。动作适合持续一段时间、需要反馈且可以取消的任务。发送目标后，系统会持续规划和控制，直到成功、失败或被取消。

可以把Action理解成一次可以持续沟通的委托。普通话题更像“我把消息告诉你”，服务更像“我问一次，你回答一次”，Action则是“请完成这个任务，过程中告诉我进展，我也可以随时取消”。导航显然属于第三种情况。

一个Action包含目标、反馈和结果。目标给出位置、朝向和坐标系；反馈可用来观察剩余距离、当前恢复次数或任务进展；结果则区分成功、取消和失败。新目标还可能抢占旧目标，所以发表导航坐标后，不能确定小车一定会到达。

第一次实验时，先在RViz中完成以下闭环：设置初始位姿，确认激光与地图贴合；发送一个近距离、无遮挡的目标；观察全局路径、局部轨迹和\code{/cmd\_vel}；最后主动取消一次任务。能完成“发送—反馈—取消—停止”，才说明Action链路和底盘停止逻辑基本可用。

\booksection{8.3 代价地图（Costmap）}

\booksubsection{8.3.1 静态层、障碍层与膨胀层}

栅格地图上只有黑、白两种区域，但导航时仅知道哪里是墙还不够。小车自身有体积，传感器也会造成定位误差，所以如果小车和墙离的太近，可能会发生碰撞。为了避免这个问题，Nav2在普通地图上又叠加了一张代价地图。

全局代价地图常见三层：静态层来自保存的地图，障碍层来自实时传感器，膨胀层在障碍物周围增加代价。叠加后，规划器不仅知道墙在哪里，也知道靠墙太近会增加风险，从而会给出一条更加安全的路线。

代价地图是一张表示“通过这里有多危险”的二维栅格。致命障碍表示机器人轮廓不能与之重叠；未知区域表示尚无可靠观测；障碍附近通常按距离形成逐渐降低的代价。规划器寻找的是一条同时满足“距离最短”和“绝对安全”两个条件的路线。\cite{macenski2023survey}。

三层各自回答不同问题：

\begin{itemize}
  \item \textbf{静态层：}地图文件中原本有哪些墙和固定设施？
  \item \textbf{障碍层：}激光或点云现在看到了哪些障碍，哪些射线经过的区域可以清为空闲？
  \item \textbf{膨胀层：}机器人离障碍多近时应开始付出更高代价？
\end{itemize}

\begin{center}
\fbox{
\begin{minipage}{0.85\textwidth}
\small
\textbf{把代价地图想成地面上的颜色}

墙壁和纸箱周围是不能进入的红色区域；离障碍较近的地方是需要谨慎通过的橙色区域；离障碍较远的地方代价逐渐降低。规划器需要在这些颜色上寻找一条长度和风险都合适的路线。
\end{minipage}
}
\end{center}

值得一提的是，障碍层的标记和清除必须配合好。只标记不清除会产生残留障碍；而射线范围、障碍高度阈值或坐标系配置错误，则会导致真实障碍漏检。排错时，建议分别检查原始 /scan、障碍层中间结果和最终代价地图，便于定位问题。

\booksubsection{8.3.2 全局与局部代价地图}

为什么已经有一张全局地图，还要再维护一张局部地图？因为看完整路线和处理眼前障碍是两种不同尺度的任务。全局代价地图覆盖较大区域，用于规划整条路线；局部代价地图只覆盖小车附近，并随着小车滚动，用于及时处理新出现的人或纸箱。

实例\code{param\_mini\_mec.yaml}使用\code{map}作为全局代价地图坐标系，使用\code{odom\_combined}作为局部代价地图坐标系，并从\code{/scan}获取障碍物。这个选择正好对应前两章建立的三层TF。

全局地图需要全局一致，因此常以\code{map}为坐标系；局部地图更关心短时间内连续、平滑的运动，因此常以\code{odom\_combined}为坐标系并设置滚动窗口。若把局部地图也强行固定在存在跳变修正的全局坐标系中，控制器可能把定位修正误认为机器人瞬间移动；若全局地图不用\code{map}，规划目标又可能与保存地图不在同一参考系。

两张地图还应分别选择尺寸和更新频率。全局地图范围大，不必追求极高刷新率；局部地图范围小，却要及时反映新障碍。局部窗口太小，快速行驶时来不及制动；太大则增加计算量，并把远处暂时无关的障碍也交给控制器处理。

\booksubsection{8.3.3 机器人轮廓与安全距离}

地图里的路径只是一条细线，真实小车却有长、宽和伸出的零件。若规划器把小车当成一个点，中心虽然绕开了桌腿，车体边缘仍可能撞上去。参数\code{footprint}用多边形描述底盘轮廓。实例文件给\code{mini\_mec}配置了矩形轮廓，但正式使用前仍应测量实际小车，并把突出部件与安装误差考虑进去。

膨胀半径也不是越大越安全。太小会贴近障碍物，太大则可能把窄通道全部判为不可通行。应先测量车体，再结合定位误差和制动距离设置安全余量。

轮廓和膨胀承担的是不同职责。\code{footprint}描述机器人真正占据的几何区域，用于碰撞检查；膨胀层在障碍周围建立软代价，让路径倾向于留出余量。不能用一个虚假的超大轮廓代替全部安全设计，也不能因为设置了膨胀半径，就忽略雷达盲区、定位误差和底盘制动距离。

测量轮廓时，以\code{base\_footprint}为原点记录各顶点，包含轮子、雷达支架和可能伸出的线缆保护结构。然后在RViz中同时显示机器人轮廓与代价地图：原地旋转时轮廓不能扫过静态墙面；通过最窄通道时，致命区域之间必须仍有可供整个轮廓通过的连续空间。

\booksection{8.4 全局路径规划与局部控制}

\booksubsection{8.4.1 路径（Path）与轨迹（Trajectory）的区别}

初学导航时，最容易把“路径”和“轨迹”当成同一个词。两者都可以画成地图上的一条线，但它们回答的问题不同。

\textbf{路径（path）}主要描述机器人\textbf{经过哪些空间位置}。二维路径可以写成一串按顺序排列的位姿
\begin{equation}
P=\{\bm{p}_0,\bm{p}_1,\ldots,\bm{p}_N\},\qquad
\bm{p}_i=[x_i,y_i,\theta_i]^{\mathrm T}.
\end{equation}
它给出了从起点到终点的几何路线，却通常没有规定机器人必须在什么时刻到达每个点，也没有完整规定沿途速度和加速度。因此，同一条路径既可以慢慢走，也可以在满足安全与动力学约束时较快地走。

\textbf{轨迹（trajectory）}描述机器人状态如何\textbf{随时间演化}，可写成
\begin{equation}
\bm{x}(t)=[x(t),y(t),\theta(t),v_x(t),v_y(t),\omega_z(t),\ldots]^{\mathrm T}.
\end{equation}
轨迹不仅包含“在哪里”，还包含“什么时候到达”和“以怎样的速度运动”；视控制器和模型而定，还可能包含加速度或控制输入。把轨迹中的时间信息去掉，只看其在平面上经过的位置，可以得到相应的几何路径；反过来，仅有路径却不能唯一确定轨迹。

可以用一次上学过程理解：地图应用给出的步行线路是路径；“8:00从宿舍出发，8:05走到路口，途中速度怎样变化”则是轨迹。在Nav2中，全局规划器通常输出供机器人跟随的全局路径；控制器结合当前速度、局部障碍和运动约束，计算短时间内可执行的运动。DWB、MPPI等控制器会预测或采样带有时间推进含义的候选轨迹，再从中选择控制量。工程界有时会把局部输出也宽泛地称作“local path”，阅读文档时应根据数据是否带有时间、速度和动力学约束判断其实际含义，而不能只看名称。

\begin{center}
\fbox{
\begin{minipage}{0.85\textwidth}
\small
\textbf{快速判断}

只回答“从哪里经过”，通常是在谈路径；还回答“何时到达、以多快速度和怎样的运动状态到达”，才是在谈轨迹。全局路径可作为控制器的参考，但不是可以直接发送给底盘执行的速度命令。
\end{minipage}
}
\end{center}

\booksubsection{8.4.2 规划器与全局路径规划}

目标已经给出，第一件事就是在地图上找路。栅格搜索把地图看作相邻格点组成的图。全局规划器负责在图中搜索低代价路线。

新手可以先观察一个现象：在地图中放入障碍物后，全局路径是否绕开障碍物；改变膨胀半径后，路径是否离墙更远。这样就能看到代价地图如何影响规划结果。

Nav2中的全局规划器不只一种。NavFn、Theta*和二维A*适合在二维栅格上寻找全向路径；Hybrid-A*和State Lattice还把朝向、最小转弯半径或预先生成的运动原语纳入搜索，更适合不能随意横移或原地转向的底盘。不同的规划器适合不同型号的机器人和不同的应用场合，需要根据实际情况来选择。

麦轮小车能够横移，属于全向底盘，因此可以使用全向规划器；但能横移不代表任何方向都同样稳定。如果实物横向打滑明显，可以在控制器和速度限制中降低$v_y$，甚至让全局路径更偏向前向行驶。规划器给出的是一条可走的路线，真正能否平稳执行还取决于后面的控制器。

一些栅格路径会出现锯齿或不连续转角。Nav2可在规划与控制之间加入平滑器，减少局部不规则、改善跟踪稳定性；但平滑后的路径仍必须进行碰撞检查，因为“看起来更圆滑”可能意味着转角被切近，离障碍反而更近\cite{macenski2023survey}。

\booksubsection{8.4.3 控制器与局部轨迹生成}

全局路径画出来以后，小车还不能直接沿着这条线瞬间移动。它必须考虑当前速度、眼前障碍、底盘能否完成这个动作，再决定下一瞬间怎样走。因此，全局路径只是一串目标点，还需要控制器持续输出真正可执行的\code{Twist}。

\begin{center}
\fbox{
\begin{minipage}{0.85\textwidth}
\small
\textbf{路线和驾驶不是一回事}

手机地图告诉驾驶员开车需要经过哪几条路，却不会替驾驶员决定下一秒方向盘转多少。全局规划器相当于手机地图，局部控制器相当于驾驶员。前者看得远，后者必须不断观察车前变化。
\end{minipage}
}
\end{center}

Nav2提供了多种控制器可供我们选择，不同控制器看问题的方式不同：

\begin{itemize}
  \item \textbf{DWB：}从当前速度附近采样一组可实现速度，把它们向前模拟成候选轨迹，再用路径距离、目标距离、障碍代价等评价函数打分。DWB是DWA思想在Nav2中的可配置实现，并支持全向底盘。
  \item \textbf{RPP：}从全局路径上选择前视点，用几何控制律跟踪，并在大曲率或靠近障碍时调低速度。它计算简单、路径跟随清楚，但遇到障碍时更倾向于减速或停车，通常依赖重新规划获得绕行路线。
  \item \textbf{MPPI：}批量采样未来控制序列，用运动模型和代价函数评价预测轨迹，再选择较优控制。它能表达更丰富的前瞻行为，但对模型、评价函数、采样数量和计算频率更敏感。
\end{itemize}

观察控制器时不要只看最终\code{/cmd\_vel}。还要看全局路径是否进入局部窗口、局部代价地图是否出现障碍、候选轨迹为何被拒绝，以及进度检查器是否认为机器人“卡住”。有路径但速度始终为零，往往不是规划器问题，而是所有局部轨迹发生碰撞、速度阈值不合理或TF数据过期。

\booksubsection{8.4.4 麦克纳姆轮全向速度约束与加速度限制}

差速小车因为不具备横向平移能力，通常令\code{linear.y}=0，需要先转向再前进。麦轮小车可以使用$v_y$横移，因此麦轮小车的控制器配置不能照搬差速底盘。

至少检查以下约束：$v_x$与$v_y$最大速度、$\omega_z$最大角速度、三个方向的加速度与减速度，以及接近目标时的最小有效速度。若参数允许的速度超过底盘真实能力，控制器的预测轨迹就会与实物运动不一致。

迁移到实物时，从仿真值的较小比例开始，先调低速度和加速度。只有直行、横移和旋转都能稳定跟踪后，再提高效率。

还要检查速度平滑和底盘死区。控制器输出的小速度如果低于电机克服静摩擦所需的阈值，小车会“有命令但不动”；阈值设得太高，又会在接近目标时反复冲过终点。速度平滑器可按动态约束限制速度变化并提高输出频率，使发送给底盘的命令更连续\cite{macenski2023survey}，但它不能修复错误的轮序、里程计尺度或控制器模型。

建议依次做四个最小实验：只给$v_x$、只给$v_y$、只给$\omega_z$、最后同时给$v_x$和$v_y$。记录命令速度、里程计速度和实际位移。如果横移误差远大于前进误差，应先降低$v_y$与横向加速度，再考虑调整控制器评价权重。

\begin{center}
\fbox{
\begin{minipage}{0.88\textwidth}
\small
\textbf{本地配置中的一个待验证点}

\code{param\_mini\_mec.yaml}把MPPI运动模型设为\code{Omni}，并给出非零的\code{vy\_max}，说明配置意图是允许横移；但同一处的\code{ay\_max}为0。这个组合是否会限制横向速度变化，取决于本地Humble插件对参数的实际解释，必须通过参数读取和横移实验确认，不能仅凭文件名断言“麦轮横移已经正确配置”。此外，\code{ObstaclesCritic}中的\code{consider\_footprint}为\code{false}，而代价地图又配置了矩形轮廓；正式实物实验前应确认候选轨迹碰撞检查是否真正考虑完整车体。
\end{minipage}
}
\end{center}

\booksection{8.5 行为树与复杂任务执行}

\booksubsection{8.5.1 行为树基础：Sequence、Fallback与Condition}

如果导航永远不会失败，任务流程只需要两步：规划，然后跟随。但真实环境中会有人挡路、地图会残留障碍、控制器也可能暂时找不到可行速度。我们需要一种结构，把“正常时怎么做”和“失败后怎么办”写清楚，这就是行为树。

序列（Sequence）表示“前一步成功后再做下一步”。回退（Fallback）表示“前一个办法失败时尝试下一个办法”。条件节点只做判断，动作节点执行任务。

\begin{center}
\fbox{
\begin{minipage}{0.85\textwidth}
\small
\textbf{用开门来理解行为树}

Sequence像“拿钥匙、开锁、推门”——前一步成功才继续。Fallback像“先刷卡，不行再输密码，还不行就找管理员”——前一个办法失败才尝试下一个。Condition则像先看一眼“门是不是已经开了”，它只判断，不负责推门。
\end{minipage}
}
\end{center}

一次简化导航可以写成：计算路径，跟随路径；若失败，则清理局部代价地图并重新计算。行为树把这种流程写成了一种可组合结构。

行为树中的每个节点在被“tick”时返回三种状态之一：成功、失败或运行中。Sequence遇到失败就停止，全部子节点成功后才成功；Fallback遇到成功就停止，只有所有办法都失败才失败；Condition应快速完成判断；调用规划器或控制器的Action则可能持续返回运行中。这个三状态接口让长时间导航任务也能被取消、抢占和监控。

行为树和普通流程图的区别在于它会周期性重新检查条件。比如“目标是否更新”“路径是否仍有效”“控制是否有进展”都可能在任务执行过程中变化。设计行为树时，需要明确每个失败由谁处理、最多重试几次，以及什么时候必须把失败返回给上层。

\booksubsection{8.5.2 Nav2行为树XML结构与任务执行器}

行为树节点会调用规划、控制和行为服务器提供的动作或服务。它负责决定何时请求规划、何时跟随、失败后做什么。

理解这一点后，排错会更清楚：没有全局路径，检查规划器和全局代价地图；有路径但小车不走，检查控制器、局部代价地图和\code{/cmd\_vel}；失败后没有恢复，再检查行为树和行为服务器。

Navigation2允许在运行时加载行为树XML。下面是一段只表达结构、不对应完整Nav2端口的教学伪代码：

\begin{verbatim}
<Fallback>
  <Sequence>
    <ComputePathToPose/>
    <FollowPath/>
  </Sequence>
  <Sequence>
    <ClearLocalCostmap/>
    <ComputePathToPose/>
    <FollowPath/>
  </Sequence>
</Fallback>
\end{verbatim}

它表达的是“先正常规划和跟随；失败后清理局部环境，再尝试一次”。真实Nav2行为树还会加入重规划频率、目标更新、路径有效性、进度检查、局部恢复和系统级恢复。Navigation2论文展示的分层恢复思想是：具体服务器失败时先执行针对性恢复，仍然失败时再逐步升级为清图、旋转或等待等系统恢复\cite{macenski2020marathon}。

\booksubsection{8.5.3 多目标点巡航：NavigateThroughPoses}

单目标导航解决的是“去一个地方”，巡检和配送却经常要求“小车依次去多个地方，并在每个地方做一件事”。这时不能只把几个坐标随手连续发布，还要定义某个点失败后是重试、跳过还是终止。

项目中的\code{nav2\_waypoint\_cycle/waypoint\_cycle.py}订阅\code{/clicked\_point}，把目标发布到\code{/goal\_pose}，并监听\code{navigate\_to\_pose/\_action/status}。到达一个目标后，它继续发布下一个目标；失败时会重试一次，再决定是否跳到下一点。

这是一个容易理解的巡航示例，但它直接读取状态数组最后一个元素，并通过目标话题驱动单目标导航。正式工程还应处理空状态列表、目标取消、节点重启和任务持久化。

Nav2还提供\code{NavigateThroughPoses}和Waypoint Follower一类多点任务接口。前者把一组位姿作为同一个导航任务处理；Waypoint Follower则依次访问航点，并可在每个点调用任务执行插件，例如等待、拍照或请求人工确认。航点序列可以循环执行，适合巡检和重复运输\cite{macenski2023survey}。

选择接口时先问“一个点失败后怎么办”。如果所有点共同组成一条不可拆分的路线，可以把它们作为整体任务；如果每个点还有独立操作，并允许记录、重试或跳过，航点跟随更容易表达。无论采用哪种方式，都应记录当前目标编号、重试次数、取消原因和最终结果。

\booksection{8.6 动态避障与恢复行为}

\booksubsection{8.6.1 动态障碍物感知与局部代价地图实时更新}

地图保存时教室里没有人，但真正导航时，学生可能突然从小车前方经过。静态地图无法提前知道这件事，所以局部代价地图必须不断接收新的传感器观测。

激光雷达看到新的障碍物后，局部代价地图在对应位置标记高代价。控制器重新评价候选速度：若还有安全轨迹，就绕行；若暂时没有，就减速或停止；若障碍持续存在，系统可能请求重新规划。

这条链中任何一环延迟过大，都可能让小车反应迟缓。因此动态避障不仅是算法问题，还依赖扫描频率、时间戳、代价地图更新频率和底盘制动能力。

“动态避障”也不等于系统准确识别出了“这是一个正在走路的人”。在基础配置中，激光观测通常只是把当前占据区域写入代价地图，控制器根据更新后的风险重新选择轨迹。只有在额外接入目标检测、跟踪或预测模块时，系统才可能显式利用障碍物速度和未来位置。教材实验应如实区分“响应实时障碍”与“预测动态目标”。


\booksubsection{8.6.2 减速、停车与重新规划}

在仿真画面中，速度可以看起来从某个值立刻变成零；真实小车却做不到。电机、轮胎和车体都有惯性，控制器发布零速度后，小车仍可能滑行一小段距离。安全距离至少应覆盖感知延迟、计算延迟、通信延迟和制动距离。

可以用一个保守关系帮助理解安全余量：

\begin{equation}
d_{\mathrm{safe}} \geq v\,t_{\mathrm{delay}} + d_{\mathrm{brake}} + d_{\mathrm{localization}},
\end{equation}

其中$t_{\mathrm{delay}}$包含传感器、代价地图、控制计算和通信延迟，$d_{\mathrm{brake}}$是当前速度与地面条件下的制动距离，$d_{\mathrm{localization}}$表示定位及轮廓误差需要的余量。这不是Nav2内部某个参数的直接公式，而是设计实物安全距离时必须考虑的物理下限。

测试动态障碍时，先使用软质大障碍物和低速。记录障碍出现时刻、\code{/scan}变化、代价地图更新、\code{/cmd\_vel}减小以及小车实际停止的时间。没有这条时间链，就无法判断延迟发生在哪里。


\booksubsection{8.6.3 恢复行为及失败处理}

小车停下来以后，下一步不应该永远是“再试一次”。它要先判断失败是否可能恢复。常见恢复包括等待、清除代价地图、原地旋转或后退。每个动作都应有适用条件。例如传感器短暂遮挡时可以等待；代价地图残留虚假障碍时可以清图；狭窄空间中盲目旋转或后退则可能更危险。

\begin{center}
\fbox{
\begin{minipage}{0.85\textwidth}
\small
\textbf{恢复不是“万能重启键”}

清图只适合处理环境模型中的陈旧或错误信息，等待只适合可能自行消失的阻塞，旋转和后退则必须先有足够安全空间。如果真正的问题是定位错误、TF中断或电机失效，反复恢复只会浪费时间，甚至带来新的风险。
\end{minipage}
}
\end{center}

恢复失败后，系统应停止并报告原因。行为树需要限制重试次数，并把失败状态返回给任务调用者。

恢复动作应与证据匹配：

\begin{itemize}
  \item 全局规划器因陈旧障碍找不到路径，可以尝试清理全局代价地图后重规划；
  \item 控制器附近出现短时行人，可以等待，或让全局规划器寻找新路线；
  \item 机器人局部被障碍包围且空间足够，可以清理局部地图后谨慎旋转；
  \item 定位已经错误、TF中断或底盘不能执行速度时，清图和旋转通常不会解决根因，应终止任务并报告故障。
\end{itemize}

在人员密集场景中，执行恢复并不一定表示性能差。真正需要警惕的是同一恢复反复出现却没有改善，或者恢复动作本身进入不安全空间。

\booksection{8.7 真实机器人导航调参与系统集成}

\booksubsection{8.7.1 仿真参数向实物迁移与初始值设置}

前面已经在算法层面把导航链路接通了。现在要把它放到真实小车上。迁移时最危险的想法是：“仿真已经成功，把同一个YAML复制过去就行。”仿真里的轮胎不会磨损，地面摩擦通常比较理想，传感器也很少突然延迟；这些条件到了实物上都会改变。

实例入口\code{wheeltec\_robot\_nav2/launch/wheeltec\_nav2.launch.py}默认参数文件在源码中指向\code{param\_mini\_akm.yaml}，而我们的主线车型是\code{mini\_mec}。因此启动时应显式传入麦轮参数文件，或在经过版本管理的教学分支中修改默认值。

\begin{verbatim}
# ROS 2 Humble；实物mini_mec；路径需按安装环境确认；待运行验证
ros2 launch wheeltec_nav2 wheeltec_nav2.launch.py \
  map:=/absolute/path/to/classroom.yaml \
  params:=/absolute/path/to/param_mini_mec.yaml \
  use_sim_time:=false
\end{verbatim}

仿真中应使用\code{use\_sim\_time:=true}，实物中使用\code{false}。如果部分节点使用仿真时钟、部分节点使用系统时钟，TF和传感器消息会出现难以理解的“过去数据”或超时错误。

不要把仿真参数文件原样复制到实物。迁移前至少核对下面六类量：

\begin{enumerate}
  \item 坐标系名称、雷达外参和地图文件路径；
  \item 机器人真实轮廓、传感器最小与最大有效距离；
  \item 最大速度、加速度、减速度和底盘最小有效速度；
  \item 控制器是否允许$v_y$，运动模型是否与麦轮底盘一致；
  \item 局部窗口尺寸、更新频率和机器人最大速度下的制动距离；
  \item 行为树文件、规划器与控制器插件在Humble环境中是否实际存在。
\end{enumerate}

第一次启动时让车轮离地或关闭电机使能，先确认\code{/cmd\_vel}方向和急停链路；落地后把速度限制到低值，只发送近距离目标。仿真成功只能证明接口和逻辑大致连通，不能证明轮胎摩擦、载荷、无线网络和制动距离满足实物要求。

\booksubsection{8.7.2 定位、代价地图与控制器串级调参}

Nav2参数很多，新手最容易犯的错误是同时打开多个YAML区域，看到哪里像问题就改哪里。更可靠的办法是沿数据流逐层验收：上一层没有通过，就不进入下一层。推荐顺序如下：

\begin{enumerate}
  \item 定位：激光与地图是否贴合，AMCL是否稳定；
  \item 全局代价地图：静态地图、障碍和膨胀是否合理；
  \item 全局规划：能否得到可行路径；
  \item 局部代价地图：新障碍能否及时出现并清除；
  \item 控制器：低速下能否沿路径行驶并停在目标附近；
  \item 行为树与恢复：故障时是否按预期重试并安全结束。
\end{enumerate}

每次只改一组相关参数，并保存运行记录。否则，结果变好或变坏时都不知道是谁造成的。

每一级都要设置清楚的“通过条件”：

\begin{itemize}
  \item 定位通过：静止时激光与地图边界贴合，短距离运动后能够稳定回到正确位置；
  \item 代价地图通过：静态墙、临时纸箱和障碍清除过程都能在RViz中对应观察；
  \item 规划通过：多个可达目标均能生成不穿越致命障碍的路径，不可达目标能明确失败；
  \item 控制通过：低速跟踪时不持续振荡、不擦碰膨胀区，并能在容差内停车；
  \item 恢复通过：人为阻塞后按预期等待、重规划或终止，不无限重试。
\end{itemize}

若某一级没有通过，不要通过调高下一层的容差掩盖问题。例如定位持续跳变时，放宽目标容差只能让系统更容易宣布“到达”，并不能让实际位置更准确。

\booksubsection{8.7.3 导航性能评价}

小车终于到达目标，并不代表调试已经结束。一次成功可能只是路线简单、现场刚好没人，也可能是某个错误没有在这次实验中暴露。要判断系统是否真的变好，必须重复实验并记录指标。

至少记录成功率、实际路径长度和任务时间。成功率说明系统是否可靠；路径长度反映绕行与规划效率；任务时间还包含减速、等待和恢复代价。

设共执行$N$次任务，其中成功$N_s$次，则成功率为

\begin{equation}
R_s=\frac{N_s}{N}\times 100\%.
\end{equation}

若第$i$次成功任务的实际路径长度和任务时间分别为$L_i$与$T_i$，可以报告成功任务的均值，同时单独统计失败任务停在哪一阶段。只对成功样本求平均不能替代失败率；一套系统可能“成功时很快”，却经常无法到达。

还应记录路径跟踪误差、最小障碍距离、恢复次数和最终位姿误差。对于动态障碍实验，增加从障碍进入扫描范围到速度首次下降的响应时间。不同算法比较必须使用相同地图、起终点、机器人限制和障碍脚本，否则差异可能来自实验条件而不是算法。

不要只展示一次成功视频。选择多个目标点，重复若干次，并记录失败类型。实物实验还应注明电量、载荷、地面和动态障碍条件。

\booksubsection{8.7.4 典型故障}

当导航失败时，终端里可能同时出现规划、控制和TF警告。不要被最后一条红色日志牵着走。故障往往从底层开始，却在更上层才表现出来。仍然按从底层到上层的顺序检查：

\begin{enumerate}
  \item 轮序、速度方向和底盘是否能及时停止；
  \item 编码器、IMU与\code{/odom\_combined}是否稳定；
  \item \code{/scan}、雷达外参与TF是否正确；
  \item 地图、初始位姿和AMCL是否收敛；
  \item 轮廓、膨胀和速度加速度是否符合实物；
  \item 规划器、控制器与行为树是否报告具体错误。
\end{enumerate}

这个顺序体现了全文主线：导航失败不一定是导航算法错误。上层系统建立在第6章和第7章之上。

可以把常见症状进一步对应到证据：

\begin{itemize}
  \item \textbf{目标发送后完全无反应：}检查Action服务器、生命周期状态和目标坐标系；
  \item \textbf{没有全局路径：}检查机器人是否位于自由空间、目标是否可达、全局地图和规划器错误码；
  \item \textbf{有路径但速度为零：}检查局部代价地图、控制器轨迹、进度检查器和速度阈值；
  \item \textbf{能走但左右振荡：}检查定位噪声、路径平滑、控制频率、加速度限制和评价权重；
  \item \textbf{接近目标反复越过：}检查最小有效速度、减速度、目标容差和底盘实际制动；
  \item \textbf{仿真成功而实物擦碰：}重新测量轮廓、外参、定位误差和制动距离，不能只增大膨胀半径碰运气。
\end{itemize}

\booksection{8.8 本章小结与习题}

本章把第6章的局部运动、第7章建立的地图与本章首先完成的AMCL定位组合成了一条完整导航链。地图加载后，系统先接收初始位姿并验证定位收敛，目标随后才进入行为树导航器；规划服务器在全局代价地图上生成路径；控制服务器结合局部代价地图和机器人运动约束生成可执行的局部运动并输出速度；遇到临时阻塞或局部失败时，行为树选择有条件、有限次的恢复动作。

需要始终记住四个边界：全局规划器负责“路线”，不直接控制电机；控制器负责短时运动，却依赖正确的定位、TF和障碍表示；代价地图表达风险，不等于完整理解动态物体；恢复行为帮助系统摆脱可恢复故障，不能修复错误坐标系、失效传感器或机械故障。

对麦轮小车而言，$v_y$让它能够横移，但也增加了模型与实物不一致的可能。导航参数必须以直行、横移、旋转三个底层实验为基础，再依次验证定位、代价地图、规划、控制和恢复。最终结果不能只是一张“到达目标”的截图，而应是一组能解释成功和失败的观测证据。

\booksubsection{8.8.1 Nav2架构与行为树理解题}
\begin{enumerate}
  \item 画出目标任务、行为树导航器、规划器服务器、控制器服务器、行为服务器、代价地图和底盘之间的数据流，并标出\code{NavigateToPose}与\code{/cmd\_vel}。
  \item 为什么已经生成全局路径时，控制器仍可能拒绝运动？
  \item 行为树中的Sequence、Fallback与Condition各适合表达哪类导航逻辑？
  \item 生命周期节点的进程已经存在，为什么仍可能不能接收任务？配置与激活分别完成什么工作？
  \item 比较NavFn一类全向栅格规划器与Hybrid-A*一类运动学可行规划器。麦轮小车在什么情况下仍可能不希望大量横移？
  \item 比较DWB、RPP和MPPI的基本思想。为什么不能只根据论文中的性能描述直接决定实物控制器？
\end{enumerate}

\booksubsection{8.8.2 多目标巡航与动态避障任务}
在仿真和实物中完成同一个任务：小车从起点出发，依次到达两个目标点，途中放置一个可移动障碍物，最后回到起点附近。分别用单目标Action串联、\code{NavigateThroughPoses}或Waypoint Follower中的一种实现，并说明选择理由。

记录障碍出现后\code{/scan}、局部代价地图、局部轨迹和\code{/cmd\_vel}的变化，测量从首次观测到首次减速的时间。若一个目标失败，说明系统采用重试、跳过还是终止；若执行恢复，记录恢复类型、次数和恢复后是否真正改善。

\booksubsection{8.8.3 真实导航调参与评价题}
选择至少三个目标点，在仿真和实物中分别重复执行导航任务。统计成功率、平均路径长度和平均任务时间，并记录一种典型失败。解释麦轮底盘的$v_y$能力、速度与加速度限制、机器人轮廓、膨胀半径和定位误差分别如何影响结果。

增加一组单变量对照实验，可从膨胀半径、最大横向速度、控制器频率或局部窗口尺寸中选择一个参数。每次只改变这一项，比较成功率、路径长度、任务时间、最小障碍距离和恢复次数，并解释变化机制。

最低提交物包括接口清单与TF主链、地图与车型参数、启动命令、行为树或任务接口说明、任务记录、代价地图和路径截图、一次故障诊断过程，以及仿真和实物结果差异分析。任何尚未运行的命令或预期现象都必须标记为“待验证”。

\booksubsection{8.8.4 综合故障诊断题}

下面各题先写观察证据，再写原因假设和检查顺序。

\begin{enumerate}
  \item RViz中能看到完整全局路径，但小车持续输出零速度。怎样判断是局部地图封路、控制器拒绝全部轨迹、进度检查失败，还是底盘没有执行命令？
  \item 行人离开后，局部代价地图仍留下障碍。怎样分别检查传感器射线、清除范围、坐标系和观测持久时间？
  \item 小车在仿真中能顺利横移绕障，实物却横向振荡。请从$v_y$里程计、轮胎打滑、加速度限制、控制频率和机器人轮廓安排检查顺序。
  \item 机器人到达目标附近却不断前后摆动。怎样区分目标容差过严、最小速度过高、定位噪声和制动不足？
  \item 行为树不断执行清图和旋转，却始终无法成功。什么证据表明故障已经超出恢复行为的适用范围，应当终止任务？
\end{enumerate}

\paragraph{Technical English}
阅读并解释下面三句话：

\begin{quote}
``The global planner computes a route, while the controller generates locally feasible velocity commands.''

``A layered costmap combines static knowledge with live sensor observations.''

``Recovery behavior should mitigate a specific failure rather than hide its root cause.''
\end{quote}

请指出\emph{route}、\emph{locally feasible}、\emph{layered costmap}、\emph{live observation}和\emph{root cause}分别对应本章哪个组件或可观察现象，并用英文写一段不少于100词的导航实验报告。报告至少包含目标、配置、结果、一次失败证据和下一步改进。

% ============================================================
\bookchapter{第9章 相机、深度感知与YOLO26视觉识别}

\section*{本章任务：让机器人看见目标，并给出目标在哪里}

前几章解决了“机器人怎样移动”和“机器人怎样到达目标点”。但是，导航地图中的一个坐标并不会自动告诉机器人，桌面上哪个物体是水杯、通道里出现的是纸箱还是行人。要让机器人根据物体类别执行观察、靠近、抓取或报告任务，还需要把相机图像变成带有语义和空间位置的结果。

本章以麦克纳姆轮机器人搭载的彩色与深度相机为主线，完成一条完整视觉链路：相机发布图像及标定参数，YOLO26从彩色图像中检测目标，深度数据把二维检测框提升为三维坐标，TF再把相机坐标变换到机器人或地图坐标系。完成本章后，读者应能够：

\begin{itemize}
  \item 区分彩色相机、双目相机与深度相机提供的信息，画出视觉感知数据链；
  \item 读取\code{Image}、\code{CameraInfo}和\code{PointCloud2}消息，检查时间戳、编码和坐标系；
  \item 用针孔模型解释内参、畸变、深度和三维坐标之间的关系；
  \item 设计相机标定实验，并用重投影误差和独立图像判断标定质量；
  \item 为机器人场景建立类别、标注和训练集划分规范，理解Precision、Recall和mAP；
  \item 设计YOLO26推理节点的输入、输出、QoS与性能观测接口；
  \item 把二维检测、对齐深度和TF连接起来，得到目标在机器人坐标系中的三维位置；
  \item 根据原始证据诊断“无图像”“检测错误”和“三维坐标错误”等视觉故障。
\end{itemize}

\noindent\statusstatic 本章以ROS~2 Humble的\code{sensor\_msgs}与\code{vision\_msgs}接口、Ultralytics YOLO26官方模型说明及项目视觉设备边界为依据。项目尚未提供确定的相机型号、驱动包、话题名称和实测标定文件，因此命令中的话题均为教学接口示例，训练速度、推理帧率、深度精度和实物效果必须在目标平台复测。YOLO26代码、模型与权重受其原许可证约束，本教材不把第三方权重作为原创配套资源重新分发。

\booksection{9.1 机器人视觉系统与任务链路}

\booksubsection{9.1.1 图像感知在机器人中的作用}

相机首先提供的是规则排列的像素，而不是“杯子”“门”或“三维位置”。视觉算法需要逐层回答三个问题：图像中出现了什么，目标位于图像哪里，目标相对机器人又在哪里。分类只回答整幅图像或某个图像块属于哪一类；目标检测还给出二维位置；深度感知和坐标变换进一步把目标放入三维空间。

因此，视觉识别并不等于机器人已经理解场景。一个检测框只能说明模型在当前图像中找到了与某类训练样本相似的区域。光照变化、遮挡、运动模糊和训练数据偏差都会改变结果。机器人在采取运动动作前，还要检查结果是否足够新、置信度是否达到任务要求、目标是否连续出现，以及三维位置是否落在合理空间内。

\booksubsection{9.1.2 彩色相机、双目相机与深度相机}

彩色相机记录光强和颜色，适合纹理、类别与外观识别，却不能从单帧图像直接唯一确定真实距离。同一个杯子离相机越远，在图像中通常越小；如果不知道杯子真实尺寸，仅凭像素宽度无法区分“小物体在近处”和“大物体在远处”。

双目相机用两个相隔一定基线的相机观察同一场景。目标在左右图像中的水平位置差称为视差。在完成标定与校正的理想模型中，深度近似满足
\begin{equation}
Z=\frac{fB}{d},
\end{equation}
其中$f$为以像素表示的焦距，$B$为双目基线，$d$为视差。视差很小时，微小匹配误差会造成较大的远距离深度误差；无纹理墙面、重复纹理和遮挡边界也会使匹配困难。

深度相机直接发布逐像素距离或点云，其测距原理可以是结构光、飞行时间或设备内部完成的双目计算。它让二维识别结果更容易变成三维坐标，但“有深度话题”不代表每个像素都有可靠数值。反光、透明、过暗、超出量程和物体边缘都可能产生零值、非数或离群点。

\booksubsection{9.1.3 从图像输入到机器人决策}

本章采用下面的数据链：

\begin{center}
彩色图像 $+$ 相机内参 $\rightarrow$ YOLO26二维检测\\
$\downarrow$\\
对齐深度 $+$ 检测区域 $\rightarrow$ 相机坐标中的三维目标\\
$\downarrow$\\
TF变换 $\rightarrow$ \code{base\_footprint}或\code{map}中的目标 $\rightarrow$ 任务决策
\end{center}

这条链中的任何一步都不能用后面的“结果看起来合理”替代检查。彩色和深度图像必须来自相符的时刻；深度必须与彩色像素对齐；内参必须对应当前分辨率；TF必须能在图像采集时刻查询。最终发布目标时，还应保留源图像时间戳、类别、置信度和坐标系，便于下游判断结果能否使用。

\booksection{9.2 ROS图像与相机接口}

\booksubsection{9.2.1 Image消息：像素之外还要看什么}

ROS~2通常用\code{sensor\_msgs/msg/Image}传递未压缩图像。\code{height}和\code{width}表示行列数，\code{encoding}说明像素含义与排列方式，\code{step}表示一行数据占用的字节数，\code{data}保存实际像素。\code{header.stamp}应对应采集时刻，\code{header.frame\_id}应指向相机光学坐标系\cite{ros2humbleimage}。

常见彩色编码有\code{rgb8}和\code{bgr8}，常见深度编码有\code{16UC1}和\code{32FC1}。不能看到三通道数组就默认它是RGB，也不能看到整数深度就默认单位是米。驱动和桥接程序必须明确编码、单位以及无效值约定。若把BGR当RGB，红蓝类别可能互换；若把毫米当米，三维坐标会放大一千倍。

接口检查可从消息元数据开始：
\begin{verbatim}
# ROS 2 Humble；话题名按实际相机驱动替换；待运行验证
ros2 topic list | grep image
ros2 topic info /camera/color/image_raw --verbose
ros2 topic hz /camera/color/image_raw
ros2 topic echo /camera/color/image_raw --once
\end{verbatim}

不要在终端持续打印完整图像数组。检查一次消息头、尺寸和编码后，应使用专用图像查看节点或编写统计节点观察亮度、丢帧和时间戳；正式教材插图使用原创示意图或实验照片，不直接使用软件界面截图。

\booksubsection{9.2.2 CameraInfo消息：把像素与光线联系起来}

\code{sensor\_msgs/msg/CameraInfo}描述图像对应的相机模型。矩阵$K$保存原始图像的内参，$D$保存畸变系数，$R$用于立体校正，$P$是校正图像的投影矩阵。其\code{width}和\code{height}应与所配套的图像分辨率一致，时间戳和\code{frame\_id}也应与图像形成清楚的对应关系\cite{ros2humbleimage}。

如果$K$的第一个元素为零，ROS客户端通常可以把相机视为尚未标定。更隐蔽的问题是加载了“另一分辨率或另一台相机”的标定文件：消息不为空，算法也能运行，但投影位置会产生系统误差。检查标定信息时不能只问“有没有\code{camera\_info}”，还要核对设备、分辨率、光学坐标系和标定日期。

\booksubsection{9.2.3 图像传输、同步与QoS}

未压缩图像数据量很大。分辨率为$W\times H$、每像素$b$字节、频率为$f$时，理想情况下单路数据率约为
\begin{equation}
R=W H b f.
\end{equation}
例如提高分辨率或帧率，会同时增加通信、内存复制和推理负担。\code{image\_transport}可以在ROS图像接口之上选择原始或压缩等传输方式，但压缩会增加编码、解码延迟，并可能损失细节。

传感器数据通常重视“尽快处理最新帧”，不希望旧图像在队列中越积越多。QoS应与相机发布者兼容，队列深度要结合处理速度设置。可靠传输并不一定适合所有高频图像链路；如果推理速度低于输入频率，盲目增大队列只会让机器人处理越来越旧的场景。工程上应同时记录输入频率、处理频率、丢帧和端到端延迟。

彩色图、深度图和\code{CameraInfo}还需要时间关联。同步不能只比较消息到达回调的时间，应比较采集时间戳。无法做到严格同步时可以使用有限容差的近似同步，但必须记录容差，并验证快速运动时的空间误差是否可接受。

\booksection{9.3 相机模型、内参与畸变}

\booksubsection{9.3.1 针孔相机模型}

针孔模型把相机坐标系中的三维点$(X,Y,Z)$投影到图像像素$(u,v)$：
\begin{equation}
u=f_x\frac{X}{Z}+c_x,\qquad
v=f_y\frac{Y}{Z}+c_y.
\end{equation}
$f_x$和$f_y$是以像素为单位的横纵焦距，$(c_x,c_y)$是主点。矩阵形式为
\begin{equation}
s
\begin{bmatrix}u\\v\\1\end{bmatrix}
=
\begin{bmatrix}
f_x&0&c_x\\0&f_y&c_y\\0&0&1
\end{bmatrix}
\begin{bmatrix}X\\Y\\Z\end{bmatrix},
\end{equation}
其中$s$与深度有关。这个模型说明：像素坐标不是三维位置，只有再知道深度，才能沿相应成像光线恢复空间点。

ROS相机光学坐标系通常规定$x$向图像右方、$y$向下、$z$指向相机前方。这与移动机器人常用的\code{base\_footprint}坐标方向不同，因此即使三维数值正确，不经过TF也不能把“相机前方一米”直接写成机器人坐标。

\booksubsection{9.3.2 焦距、主点与镜头畸变}

真实镜头不会完全服从理想针孔模型。径向畸变会让直线在图像边缘向外或向内弯曲，切向畸变常与镜头和成像平面未完全对正有关。常用模型用$k_1,k_2,k_3$等描述径向项，用$p_1,p_2$描述切向项。

畸变最容易在画面边缘暴露。若目标检测只关心类别，轻微畸变可能暂时不明显；但当检测框中心要与深度图对应，或要把像素反投影为三维点时，畸变与校正状态必须统一。不能用原始畸变图像的像素坐标配合校正图像的投影矩阵，也不能把一组内参同时套用到任意裁剪、缩放后的图像上。

\booksubsection{9.3.3 标定文件与分辨率一致性}

相机标定文件通常记录图像尺寸、相机矩阵、畸变模型及系数、校正矩阵和投影矩阵。文件路径只是配置入口，真正运行时应以发布出来的\code{CameraInfo}为准，检查驱动是否确实加载了预期文件。

图像缩放后，焦距与主点应按相同比例调整；裁剪后，主点还会随裁剪原点改变。最安全的做法是让驱动针对实际输出模式发布匹配的\code{CameraInfo}，而不是在推理节点中凭经验修改几个数字。若算法为了YOLO输入而把图像缩放或填充到方形尺寸，还应保留从网络输入坐标映射回原图坐标的比例和边界填充值。

\booksection{9.4 相机标定与质量验证}

\booksubsection{9.4.1 标定板与数据采集}

常见标定目标包括棋盘格、圆点阵列和带编码标记的标定板。关键不在于“拍够某个固定张数”，而在于让角点覆盖视野中央与四周，同时包含不同距离、倾斜角和方向。若所有图像都正对标定板且只出现在画面中心，算法很难可靠估计边缘畸变和部分参数耦合。

采集前应固定相机焦距、分辨率和成像模式，保证标定板平整，并准确测量格长。运动模糊、过曝、反光以及角点只占很小区域的图像应剔除。双目标定还要求左右相机同步观察同一块板，并保持相机之间的机械结构不变。

\booksubsection{9.4.2 内参与畸变参数计算}

标定算法根据已知标定板几何和检测到的图像角点，求解使投影误差较小的内参与畸变参数。ROS~2可使用\code{camera\_calibration}等工具完成采集和求解，但具体命令受安装版本、板型和话题名称影响，应在教学平台确认后给出。

一次标定的输出不是“永远正确”的设备属性。改变分辨率、数字裁剪、对焦状态或镜头安装，都可能使原参数失效。把标定结果写入文件后，应记录相机序列号或设备标识、图像模式、板格尺寸、日期和操作者，以避免文件被错误复用。

\booksubsection{9.4.3 重投影误差与标定质量}

设第$i$个角点的实测像素为$\bm{q}_i$，根据所求参数重新投影得到$\hat{\bm{q}}_i$，均方根重投影误差可写成
\begin{equation}
e_{\mathrm{rms}}=\sqrt{\frac{1}{N}\sum_{i=1}^{N}\left\|\bm{q}_i-\hat{\bm{q}}_i\right\|^2}.
\end{equation}
数值越小通常越好，但单个平均值不能说明全部问题。大量相似视角可能得到很小的训练内误差，却在画面边缘或新距离上表现很差。

质量判断至少包含三类证据：查看校正后直线是否仍明显弯曲；用未参与求解的独立图像计算误差；在已知尺寸或距离的实物上检查三维恢复结果。若某些视角误差特别大，应先检查角点、板尺寸和图像质量，而不是机械地删除所有“不好看”的样本来降低平均值。

\booksection{9.5 深度图像、反投影与点云}

\booksubsection{9.5.1 深度值与无效测量}

深度图的每个像素保存一个距离相关数值，但必须先确认它表示沿光轴的$Z$，还是沿视线方向的量程。常见\code{16UC1}深度可能以毫米表示，\code{32FC1}常以米表示，但最终应以设备驱动文档和消息约定为准，不能只根据编码猜单位。

无效深度可能表示为0、NaN或特殊上限值。检测框中心恰好落在反光标签、孔洞或物体边缘时，只读取一个像素非常脆弱。更稳健的方法是在检测框内部选择中央区域，过滤无效值后取中位数或稳健统计量，同时保留有效像素比例。区域也不能过大，否则背景深度会混入小目标结果。

\booksubsection{9.5.2 从像素坐标恢复三维坐标}

对于已经校正、与深度对齐的图像，已知像素$(u,v)$处的光轴深度$Z$，可由
\begin{equation}
X=\frac{(u-c_x)Z}{f_x},\qquad
Y=\frac{(v-c_y)Z}{f_y},\qquad
Z=Z
\end{equation}
得到相机光学坐标系中的三维点。这个公式很简单，前提却很多：像素必须属于与内参匹配的图像；深度必须与该像素对应；单位必须统一；畸变状态必须一致。

误差还会随距离传播。若深度误差为$\delta Z$，像素和内参也存在误差，则$X$和$Y$会同时受到影响。目标越远、越靠画面边缘，标定与像素定位误差往往越明显。因此三维定位结果应报告合理精度范围，不能因为输出保留了很多小数位就声称具有毫米级准确度。

\booksubsection{9.5.3 PointCloud2数据观察}

\code{sensor\_msgs/msg/PointCloud2}把大量点组织在一个二进制缓冲区中。\code{fields}描述每个点包含哪些字段及其偏移，常见字段为$x$、$y$、$z$，还可能包含颜色、强度或其他属性。\code{point\_step}表示一个点占用的字节数，\code{row\_step}表示一行占用的字节数，\code{is\_dense}说明是否保证没有无效点\cite{ros2humbleimage}。

有组织点云保留与图像相似的行列结构，便于从像素查找对应点；无组织点云更像一维点列表。处理前应检查点云的\code{frame\_id}、频率、字段布局和非有限值比例。不要把点云中数组下标相同理解成它必然与另一幅彩色图像像素对齐，只有驱动明确发布了对齐数据并且分辨率一致时，这种对应才成立。

\booksection{9.6 YOLO26目标检测基础}

\booksubsection{9.6.1 检测任务的输入与输出}

目标检测以图像为输入，输出若干目标候选。每个结果通常包含类别、置信度和边界框；一张图可以没有目标，也可以有多个同类目标。检测不同于图像分类，因为它还要回答目标在图像哪里；检测也不同于实例分割，因为矩形框不精确描述物体轮廓。

Ultralytics YOLO26提供不同规模的检测模型，并支持训练、验证、推理和导出等工作模式\cite{ultralytics2026yolo26}。模型规模越大通常需要更多计算与显存，但具体精度和速度必须在目标数据与硬件上测量。教材以较小模型建立基线，不把官方基准直接当作机器人平台的实测性能。

\booksubsection{9.6.2 类别、边界框与置信度}

边界框可以用两角点$(x_1,y_1,x_2,y_2)$表示，也可以用中心、宽和高$(c_x,c_y,w,h)$表示；有的训练格式还会把坐标归一化到$[0,1]$。接口转换时必须明确格式、坐标原点、是否归一化，以及坐标对应原图还是经过缩放和填充的网络输入。

置信度是模型对当前预测的评分，不是“目标存在的客观概率”，也不是跨模型、跨类别通用的质量证书。阈值过高会增加漏检，阈值过低会增加误检。机器人任务应根据后果选择阈值：用于提示操作者时可以容忍更多候选；用于触发靠近或抓取时，通常还要结合连续帧、深度范围和任务状态进行复核。

交并比（Intersection over Union，IoU）衡量两个框的重叠程度：
\begin{equation}
\mathrm{IoU}=\frac{|B_{\mathrm{pred}}\cap B_{\mathrm{gt}}|}{|B_{\mathrm{pred}}\cup B_{\mathrm{gt}}|}.
\end{equation}
它用于判断预测框与标注框是否匹配，也是计算检测评价指标的重要基础。

\booksubsection{9.6.3 YOLO26模型与ROS节点边界}

YOLO26负责从输入图像产生检测结果，ROS节点负责订阅、时间管理、消息转换、参数、发布和诊断。把两者分开有助于测试：模型可以先对离线测试集验证，ROS封装则可以用固定图像检查话题和时间戳。若系统输出错误，应先判断是模型预测错误，还是节点在颜色转换、缩放还原或消息发布时出错。

YOLO26默认端到端检测头与传统需要额外非极大值抑制的输出形式可能不同\cite{ultralytics2026yolo26}。若使用Ultralytics高级接口，后处理由所选模式与库完成；若导出ONNX等格式后自行解析张量，必须依据具体导出设置核对输出布局，不能照搬旧版本YOLO的后处理代码。

\booksection{9.7 机器人场景数据采集与标注}

\booksubsection{9.7.1 场景图像采集}

训练数据应覆盖机器人真正会遇到的变化：不同距离、视角、光照、背景、遮挡、运动速度和相机曝光。只在桌面上近距离拍摄清晰正面图，模型很可能在机器人移动产生模糊、目标位于画面边缘或教室光照改变时失效。

采集时应记录场景、日期、设备和采集批次。视频抽帧虽然方便，但相邻帧高度相似；若随后随机拆分数据，相邻帧可能同时进入训练集与测试集，产生严重的数据泄漏。更合理的方法是按采集序列、场景或日期分组划分。

数据采集还涉及隐私与授权。包含人脸、学生信息或受限制场所的图像，必须按学校与项目要求处理。数据可以用于课堂实验，不代表可以自动公开到教材仓库。

\booksubsection{9.7.2 类别定义与标注规范}

类别首先要服务任务。若机器人只需要发现“通道障碍物”，把纸箱按颜色拆成多个类别可能没有意义；若任务要求分别搬运不同容器，类别边界就必须清楚。标注规范应写明：遮挡到什么程度仍标注，截断目标怎样处理，反射或图片中的目标是否计入，以及多个物体重叠时分别怎样画框。

边界框应紧贴可见目标，并保持全体标注者的一致性。漏标会把真实目标当作背景，松紧不一致会增加定位噪声。同一批数据应抽样复核，统计空标注、越界框、非法类别和极小框。不要等训练失败后才第一次检查标签。

\booksubsection{9.7.3 训练集、验证集与测试集}

训练集用于更新模型参数，验证集用于选择模型与阈值，测试集只用于最终评价。测试集如果反复参与调参，就不再是独立测试。三个集合除了数量比例，更重要的是避免同源泄漏，并覆盖部署场景。

还应观察类别数量、目标尺寸和场景分布。样本总数很多，不代表每个类别和每种困难条件都足够。对机器人小目标，可以单独统计目标框相对图像面积；对深度定位任务，还应保证测试集中包含不同距离和无效深度情形。

\booksection{9.8 YOLO26训练与评价}

\booksubsection{9.8.1 训练配置与基本过程}

训练前需要确定数据配置、预训练模型、输入尺寸、批量大小、训练轮数和运行设备。一个可复现记录至少包括软件版本、模型文件来源、数据版本、随机种子、关键参数和最终权重的校验信息。下面只表示Ultralytics接口形式，不是本项目已运行记录：
\begin{verbatim}
# YOLO26接口示例；依赖、数据路径、GPU与结果均待验证
yolo detect train model=yolo26n.pt data=/absolute/path/dataset.yaml \
  epochs=100 imgsz=640
\end{verbatim}

训练轮数多不一定更好。应同时观察训练集与验证集损失和指标，保存可复现配置，并在独立测试集上评价最终模型。若计算资源有限，可以先用较小模型和小规模数据检查标签、目录和完整流程，再进行正式训练。

\booksubsection{9.8.2 Precision、Recall与mAP}

设真正例、假正例和假负例分别为$TP$、$FP$与$FN$，则
\begin{equation}
\mathrm{Precision}=\frac{TP}{TP+FP},\qquad
\mathrm{Recall}=\frac{TP}{TP+FN}.
\end{equation}
Precision高说明发布的检测中误报较少，Recall高说明真实目标较少被漏掉。改变置信度阈值通常会使二者此消彼长，因此不能只报告单一阈值下的一个数字。

平均精度（Average Precision，AP）概括某一类别在不同置信度阈值下的Precision--Recall关系，mAP再对类别取平均。报告mAP时必须说明IoU条件，例如\code{mAP@0.5}与在多个IoU阈值上平均的\code{mAP@0.5:0.95}不是同一个指标。机器人任务还应分别报告关键类别，避免总体平均值掩盖某个安全相关类别表现很差。

\booksubsection{9.8.3 过拟合、漏检与误检分析}

训练指标继续改善而验证指标停滞或下降，是过拟合的常见信号，但原因不只模型太大，还可能包括训练与验证分布不同、标注错误或数据泄漏。改进手段可以是补充真实变化、清理标签、数据增强、正则化或选择较合适模型，而不是只增加训练轮数。

漏检和误检应按条件分组分析。例如，把漏检分成小目标、遮挡、逆光和运动模糊，把误检分成相似物体、背景纹理和反射。这样的错误清单才能指导下一轮数据采集。只展示几张成功图像或只报告最高mAP，无法证明模型适合机器人运行环境。

\booksection{9.9 YOLO26推理与ROS集成}

\booksubsection{9.9.1 图像订阅与推理节点}

一个最小推理节点需要完成：订阅图像，按编码转换为模型输入，执行推理，把坐标映射回原始图像，并发布结构化检测结果。模型加载应在节点初始化阶段完成，不能每帧重新加载。阈值、模型路径、设备和输入话题应作为参数，而不是散落在回调函数中的常量。

如果推理回调耗时较长，单线程执行器可能阻塞其他回调。可以把图像接收与推理解耦，只保留最新帧，或使用工作线程和受控队列；但必须明确共享数据和节点关闭逻辑。优化前先测量预处理、模型推理、后处理和ROS发布各自耗时，不要把全部延迟都归因于神经网络。

\booksubsection{9.9.2 检测结果消息设计}

课堂演示可以定义简单消息，但工程接口更适合复用\code{vision\_msgs/msg/Detection2DArray}等标准消息\cite{ros2humblevision}。一次发布应包含源图像的\code{header}和若干检测，每个检测包含边界框与类别假设。若项目需要与不具备\code{vision\_msgs}的环境兼容，可以设计自定义消息，但至少要保留：类别标识、分数、边界框格式、图像尺寸、时间戳和坐标系。

不要只发布画好框的图像。可视化图适合人看，却迫使下游再次解析像素，丢失精确类别和置信度。结构化检测结果用于决策，调试图像作为可选输出，两者职责不同。检测消息的时间戳应继承源图像采集时间，而不是推理结束时间；另行记录处理完成时间即可计算延迟。

\booksubsection{9.9.3 推理频率、延迟与陈旧结果}

推理频率是单位时间完成多少帧，端到端延迟是从采集到结果可用经历多久，两者不能互相替代。系统可以每秒输出很多帧，却始终在处理积压的旧图像；也可能低频运行，但每次都基于最新画面。

设图像采集时间为$t_{\mathrm{capture}}$，结果发布完成时间为$t_{\mathrm{publish}}$，则端到端延迟近似为
\begin{equation}
T_{\mathrm{e2e}}=t_{\mathrm{publish}}-t_{\mathrm{capture}}.
\end{equation}
应同时记录输入频率、输出频率、平均与高分位延迟、队列长度和丢帧。机器人运动时，陈旧检测会造成明显空间偏差，因此任务层还应设置最大结果年龄，超时后停止使用该结果，而不是继续驱动机器人追逐过去的位置。

\booksection{9.10 从二维检测到三维目标定位}

\booksubsection{9.10.1 检测框与深度图对齐}

彩色传感器与深度传感器通常具有不同视点、分辨率和视场。所谓“对齐深度”，是利用两套内参与相机间外参，把深度重新投影到彩色图像像素网格。只有使用与检测图像对齐的深度，才能在检测框内直接查找距离。

对齐还包含时间条件。机器人或目标运动时，即使空间标定正确，不同时刻的彩色和深度也会错位。应检查时间戳差，并在实验中让目标经过画面边缘，观察深度轮廓是否跟随彩色边界。静止场景中的偶然重合不能证明动态同步正确。

\booksubsection{9.10.2 目标三维坐标计算}

最简单方法是在检测框中心读取深度并反投影，但中心点可能落在空洞或背景。教学实验可采用以下流程：将检测框缩小为中央区域；过滤零值、NaN和量程外数据；计算有效深度中位数；选择接近中位数的像素作为代表点；再用内参恢复$(X,Y,Z)$。

目标检测框并不等于物体三维中心。对于细长、不规则或部分遮挡物体，中央区域也可能包含背景。因此结果应称为“由检测区域估计的目标位置”，并保留有效深度比例与估计方法。若任务要求精确抓取，通常还需要实例分割、点云聚类、姿态估计或几何模型，而不是把二维框中心直接当作抓取点。

\booksubsection{9.10.3 从相机坐标变换到机器人坐标}

设相机光学坐标中的齐次点为$\tilde{\bm{p}}_c=[X_c,Y_c,Z_c,1]^{\mathrm T}$，相机到机器人基座的变换为${}^{b}\bm{T}_{c}$，则
\begin{equation}
{}^{b}\tilde{\bm{p}}={}^{b}\bm{T}_{c}\,{}^{c}\tilde{\bm{p}}.
\end{equation}
TF查询应使用图像采集时刻，而不是无条件获取“最新变换”。对于固定安装相机，外参通常为静态TF；对于安装在云台或机械臂上的相机，变换还依赖关节状态和对应时间戳。

发布三维目标时应使用\code{PointStamped}、\code{PoseStamped}或适合任务的检测消息，并明确目标坐标系。只有位置而没有姿态时，不应虚构一个“准确姿态”。如果下游导航需要\code{map}坐标，先确认AMCL定位和TF在采集时刻有效；如果只用于近距离避障或观察，\code{base\_footprint}坐标可能更直接、更不易受到全局定位修正影响。

\booksection{9.11 视觉系统的分层故障诊断}

\booksubsection{9.11.1 图像Topic无数据或图像异常}

从底层向上检查：设备是否被系统识别，驱动节点是否存活，话题名称和命名空间是否正确，发布者与订阅者QoS是否兼容，图像时间戳是否前进，编码与尺寸是否合理。节点存在但没有数据，和数据存在但订阅不到，是两类不同故障。

如果有图像但颜色错误，先检查\code{encoding}和颜色通道；如果画面间歇冻结，记录频率与时间戳而不是只看窗口；如果延迟持续增加，检查队列和处理速度。USB带宽、供电和多个高分辨率数据流也可能造成驱动层丢帧，不能通过调整YOLO阈值解决。

\booksubsection{9.11.2 彩色与深度未对齐}

先用静止、边缘清楚且距离已知的物体检查空间对齐，再让物体或机器人缓慢运动检查时间同步。若静止时也错位，重点检查分辨率、内参、彩深外参和对齐话题；若静止正确而运动时错位，重点检查时间戳、同步策略与处理延迟。

还要区分“显示缩放造成的视觉错觉”和原始数据错误。记录同一时刻彩色与深度尺寸，在若干已知像素处读取深度，并检查点云边缘。不要只凭两张经过不同窗口缩放的显示图宣布对齐失败。

\booksubsection{9.11.3 检测正确但三维坐标错误}

二维框正确只能排除一部分模型问题。三维坐标错误时按以下顺序检查：深度编码和单位，框坐标是否已映射回原图，对齐深度是否与彩色图同尺寸同时间，内参是否对应当前分辨率，反投影公式是否使用正确坐标约定，最后检查TF方向、父子坐标系和查询时间。

可以设计逐级真值：先把已知平面放在相机正前方检查$Z$；再把目标移动到画面左右检查$X$符号；上下移动检查$Y$符号；最后测量相机相对\code{base\_footprint}的安装位置，检查TF变换后的坐标。逐级测试比直接让导航系统追踪目标更容易找到根因。

\booksection{9.12 Technical English：机器人视觉}

阅读并解释下面三句话：
\begin{quote}
``Camera calibration estimates the parameters required to relate three-dimensional rays to image pixels.''

``A high-confidence detection is not automatically a geometrically valid target.''

``Depth, color, and transforms must refer to compatible frames and acquisition times.''
\end{quote}

其中，\emph{calibration}强调参数估计与验证，\emph{detection}表示从图像中找到并分类目标，\emph{depth}表示与距离相关的测量。\emph{false positive}表示系统报告了并不存在或类别错误的目标，\emph{false negative}表示真实目标被漏掉。描述实验时，应说明现象、数据条件、证据和限制，不能只写“the model works well”。

请阅读一次模型训练日志或评价表，找出Precision、Recall和mAP的取值与条件；再用英文写一段不少于100词的视觉实验报告。报告至少包含相机输入、目标类别、模型或阈值、延迟、一次漏检或误检，以及下一步改进。若三维定位尚未完成，应明确写出缺少的是深度对齐、内参还是TF证据。

\booksection{9.13 本章小结与习题}

本章建立了从像素到机器人坐标的视觉链。\code{Image}提供像素与采集时刻，\code{CameraInfo}提供投影模型，标定把理想针孔模型与真实镜头联系起来；YOLO26给出二维语义检测，对齐深度把检测区域提升到相机三维坐标，TF再把结果放入机器人任务使用的参考系。

需要始终记住四个边界：检测框不是物体精确轮廓；置信度不是几何正确性的证明；单个深度像素不一定代表目标距离；相机坐标中的正确点若使用了错误TF或时间戳，在机器人坐标中仍然会错。视觉系统应以分层证据验收，而不是只提交一张画框图片。

\booksubsection{9.13.1 相机模型与标定题}
\begin{enumerate}
  \item 已知$f_x=600$、$f_y=600$、$c_x=320$、$c_y=240$，像素$(380,210)$的光轴深度为\SI{2}{\metre}，计算其相机光学坐标。说明$x$、$y$方向。
  \item 为什么改变图像分辨率或裁剪区域后不能无条件沿用原内参？焦距和主点分别怎样受影响？
  \item 设计一组覆盖画面中心、边缘、远近与倾斜角度的标定板采集方案，并列出应剔除的图像。
  \item 重投影误差很小，是否足以证明三维测距准确？还需要哪些独立验证？
  \item 比较彩色相机、双目相机和深度相机的输出、优势及典型失效条件。
\end{enumerate}

\booksubsection{9.13.2 数据标注、训练与ROS检测任务}

选择两至三个与机器人任务有关的类别，完成一个小规模但可审计的数据闭环：制定类别与遮挡标注规范；按采集序列划分训练、验证和测试集；检查类别与目标尺寸分布；训练或微调YOLO26小型模型；在独立测试集上报告各类别Precision、Recall和mAP；最后把模型封装为ROS~2图像订阅与\code{Detection2DArray}发布节点。

至少记录输入与输出话题、消息类型、图像编码、QoS、模型及数据版本、阈值、输入频率、推理频率和端到端延迟。训练和ROS节点命令若未在Ubuntu 22.04、ROS~2 Humble和目标计算平台运行，必须标记为“待验证”。第三方权重只记录合法来源与版本，不作为原创文件重新授权。

\booksubsection{9.13.3 三维定位与视觉诊断题}

在仿真或实物中选择一个目标，发布二维检测和机器人坐标系中的三维位置。用卷尺、仿真真值或固定标定物建立参考，分别在近、中、远距离及画面中央、边缘位置重复测量。报告有效深度比例、位置误差、结果年龄和失败样本。

下面各题要求先列原始证据，再提出原因假设与检查顺序：
\begin{enumerate}
  \item 图像话题频率正常，推理节点却收不到图像。怎样检查话题名称、QoS、编码与执行器阻塞？
  \item 模型在离线图片上检测正确，机器人运动时大量漏检。怎样区分运动模糊、曝光、输入积压和场景分布差异？
  \item 二维框正确，但框中心深度经常为零。怎样设计区域深度估计，同时避免混入背景？
  \item 相机坐标中的$Z$正确，变换到\code{base\_footprint}后左右方向相反。应怎样检查光学坐标约定、TF方向和外参？
  \item 静止目标定位正确，机器人移动时目标坐标拖尾。怎样通过时间戳、同步差、端到端延迟和TF查询时刻定位原因？
  \item 模型总体mAP较高，但任务所需的小目标经常漏检。为什么总体指标会掩盖问题？下一轮应怎样补充数据和分组评价？
\end{enumerate}

最低提交物包括：视觉数据流与坐标链、相机接口清单、标定文件及质量证据、数据集版本和划分方法、训练配置与指标、结构化检测消息、三维定位误差记录、一次完整故障诊断，以及已经验证与仍待验证内容的边界。

\section*{参考资料与编写状态}

\begin{enumerate}
  \item 本地实例工程：\code{x5\_src\_250401}。本文仅依据源码结构、启动文件和参数文件进行静态审查；其中第三方代码继续受原许可证约束。
  \item 本地Navigation2源码：\code{navigation2-main}与\code{x5\_src\_250401/src/navigation2-humble}，用于核对Nav2组件、插件和动作接口。
  \item 第8章主要参考\code{references/navigation2.pdf}和\code{references/modern mobile robotics algorithms.pdf}；第7章另参考\code{references/joss.02783.pdf}及多传感器SLAM综述，用于技术背景与术语核查。
  \item 第9章的ROS消息字段以ROS~2 Humble官方接口为基准，YOLO26模式与版本特性以Ultralytics官方资料为基准；相机型号、驱动话题、标定结果、训练结果与推理性能尚待教学平台验证。
  \item YOLO26代码、模型、文档与权重继续受Ultralytics的AGPL-3.0或商业许可约束。教材仓库默认不重新分发第三方权重，正式发布前还应按项目许可策略记录来源、版本和适用许可。
  \item ROS 2主教学环境按项目决策采用Ubuntu 22.04与ROS 2 Humble；Gazebo、实物接口、命令和参数仍需逐项运行验证。
\end{enumerate}

\begin{thebibliography}{9}
\bibitem{macenski2021slamtoolbox}
Macenski, S., and Jambrecic, I.
SLAM Toolbox: SLAM for the Dynamic World.
\emph{Journal of Open Source Software}, 6(61), 2783, 2021.
\url{https://doi.org/10.21105/joss.02783}.

\bibitem{macenski2020marathon}
Macenski, S., Martín, F., White, R., and Ginés Clavero, J.
The Marathon 2: A Navigation System.
\emph{2020 IEEE/RSJ International Conference on Intelligent Robots and Systems (IROS)}, 2020.
\url{https://doi.org/10.1109/IROS45743.2020.9341207}.

\bibitem{macenski2023survey}
Macenski, S., Moore, T., Lu, D. V., Merzlyakov, A., and Ferguson, M.
From the Desks of ROS Maintainers: A Survey of Modern \& Capable Mobile Robotics Algorithms in the Robot Operating System 2.
\emph{Robotics and Autonomous Systems}, 168, 104493, 2023.
\url{https://doi.org/10.1016/j.robot.2023.104493}.

\bibitem{ros2humbleimage}
Open Robotics.
\emph{sensor\_msgs Message Definitions: Image, CameraInfo and PointCloud2, ROS 2 Humble}.
\url{https://docs.ros.org/en/humble/p/sensor_msgs/}.

\bibitem{ros2humblevision}
ROS 2 Vision Working Group.
\emph{vision\_msgs Message Definitions, ROS 2 Humble}.
\url{https://docs.ros.org/en/ros2_packages/humble/api/vision_msgs/}.

\bibitem{ultralytics2026yolo26}
Jocher, G., Qiu, J., Liu, M., Lyu, S., Akyon, F. C., and Kalfaoglu, M. E.
Ultralytics YOLO26: Unified Real-Time End-to-End Vision Models.
\emph{arXiv preprint arXiv:2606.03748}, 2026.
\url{https://doi.org/10.48550/arXiv.2606.03748}.
\end{thebibliography}

\end{document}
